
\documentclass{sebase}
%%%%%%%%%%%%%%%%%%%%%%%%%%%%%%%%%%%%%%%%%%%%%%%%%%%%%%%%%%%%%%%%%%%%%%%%%%%%%%%%%%%%%%%%%%%%%%%%%%%%%%%%%%%%%%%%%%%%%%%%%%%%%%%%%%%%%%%%%%%%%%%%%%%%%%%%%%%%%%%%%%%%%%%%%%%%%%%%%%%%%%%%%%%%%%%%%%%%%%%%%%%%%%%%%%%%%%%%%%%%%%%%%%%%%%%%%%%%%%%%%%%%%%%%%%%%
\usepackage{amsmath}
\usepackage{MIROSREPORT}

\setcounter{MaxMatrixCols}{10}
%TCIDATA{OutputFilter=LATEX.DLL}
%TCIDATA{Version=5.00.0.2606}
%TCIDATA{<META NAME="SaveForMode" CONTENT="1">}
%TCIDATA{BibliographyScheme=Manual}
%TCIDATA{Created=Thu Aug 09 10:55:09 2001}
%TCIDATA{LastRevised=Tuesday, November 22, 2005 23:55:23}
%TCIDATA{<META NAME="GraphicsSave" CONTENT="32">}
%TCIDATA{<META NAME="DocumentShell" CONTENT="Style Editor\MIROSREPORT">}
%TCIDATA{Language=American English}
%TCIDATA{CSTFile=selart.cst}

\input{tcilatex}

\begin{document}

\SetTitle{1300/TN/007 rev. 02}

%TCIMACRO{\TeXButton{frontPageMiros}{\frontPageMiros{}}}%
%BeginExpansion
\frontPageMiros{}%
%EndExpansion
\begin{tabular}{p{3cm}p{1.5cm}p{1.5cm}p{1.5cm}p{1.5cm}p{1.5cm}p{1.5cm}}
\hline\hline
\multicolumn{1}{||l}{%
%TCIMACRO{\TeXButton{Box}{\raisebox{0pt}[10mm][10mm]{}}}%
%BeginExpansion
\raisebox{0pt}[10mm][10mm]{}%
%EndExpansion
\emph{Supplier}} & \multicolumn{6}{|c||}{\FRAME{itbpF}{7.5542cm}{2.201cm}{0cm%
}{}{}{miroslogo.eps}{\special{language "Scientific Word";type
"GRAPHIC";maintain-aspect-ratio TRUE;display "USEDEF";valid_file "F";width
7.5542cm;height 2.201cm;depth 0cm;original-width 15.8823in;original-height
4.6293in;cropleft "0";croptop "1";cropright "1";cropbottom "0";filename
'MIROSLOGO.eps';file-properties "XNPEU";}}} \\ \hline
\multicolumn{1}{||l}{%
%TCIMACRO{\TeXButton{Box}{\raisebox{0pt}[5mm][2mm]{}}}%
%BeginExpansion
\raisebox{0pt}[5mm][2mm]{}%
%EndExpansion
\emph{Document No.}} & \multicolumn{6}{|l||}{1300/TN/007} \\ \hline
\multicolumn{1}{||l}{%
%TCIMACRO{\TeXButton{Box}{\raisebox{0pt}[7mm][4mm]{}}}%
%BeginExpansion
\raisebox{0pt}[7mm][4mm]{}%
%EndExpansion
\emph{Document Title}} & \multicolumn{6}{|l||}{Design of IIR Filters Using
Bilinear Transformation Method} \\ \hline
\multicolumn{1}{||l}{%
%TCIMACRO{\TeXButton{Box}{\raisebox{0pt}[5mm][2mm]{}}}%
%BeginExpansion
\raisebox{0pt}[5mm][2mm]{}%
%EndExpansion
\emph{Project}} & \multicolumn{6}{|l||}{1300-100} \\ \hline\hline
\multicolumn{1}{||l}{%
%TCIMACRO{\TeXButton{Box}{\raisebox{0pt}[5mm][2mm]{}}}%
%BeginExpansion
\raisebox{0pt}[5mm][2mm]{}%
%EndExpansion
\emph{Revision No.}} & \multicolumn{1}{|c}{\emph{1}} & \multicolumn{1}{|c}{%
\emph{2}} & \multicolumn{1}{|c}{\emph{3}} & \multicolumn{1}{|c}{\emph{4}} & 
\multicolumn{1}{|c}{\emph{5}} & \multicolumn{1}{|c||}{\emph{6}} \\ \hline
\multicolumn{1}{||l}{%
%TCIMACRO{\TeXButton{Box}{\raisebox{0pt}[5mm][2mm]{}}}%
%BeginExpansion
\raisebox{0pt}[5mm][2mm]{}%
%EndExpansion
\emph{Date}} & \multicolumn{1}{|c}{02.04.02} & \multicolumn{1}{|c}{23.05.02}
& \multicolumn{1}{|c}{} & \multicolumn{1}{|c}{} & \multicolumn{1}{|c}{} & 
\multicolumn{1}{|c||}{} \\ \hline
\multicolumn{1}{||l}{%
%TCIMACRO{\TeXButton{Box}{\raisebox{0pt}[5mm][2mm]{}}}%
%BeginExpansion
\raisebox{0pt}[5mm][2mm]{}%
%EndExpansion
\emph{Prepared By}} & \multicolumn{1}{|c}{RA} & \multicolumn{1}{|c}{RA} & 
\multicolumn{1}{|c}{} & \multicolumn{1}{|c}{} & \multicolumn{1}{|c}{} & 
\multicolumn{1}{|c||}{} \\ \hline
\multicolumn{1}{||l}{%
%TCIMACRO{\TeXButton{Box}{\raisebox{0pt}[5mm][2mm]{}}}%
%BeginExpansion
\raisebox{0pt}[5mm][2mm]{}%
%EndExpansion
\emph{Checked By}} & \multicolumn{1}{|c}{\O G} & \multicolumn{1}{|c}{\O G} & 
\multicolumn{1}{|c}{} & \multicolumn{1}{|c}{} & \multicolumn{1}{|c}{} & 
\multicolumn{1}{|c||}{} \\ \hline
\multicolumn{1}{||l}{%
%TCIMACRO{\TeXButton{Box}{\raisebox{0pt}[5mm][2mm]{}}}%
%BeginExpansion
\raisebox{0pt}[5mm][2mm]{}%
%EndExpansion
\emph{Approved By}} & \multicolumn{1}{|c}{\O G} & \multicolumn{1}{|c}{\O G}
& \multicolumn{1}{|c}{} & \multicolumn{1}{|c}{} & \multicolumn{1}{|c}{} & 
\multicolumn{1}{|c||}{} \\ \hline\hline
&  &  &  &  &  &  \\ \hline\hline
\multicolumn{7}{||l||}{%
%TCIMACRO{\TeXButton{Box}{\raisebox{0pt}[5mm][2mm]{}}}%
%BeginExpansion
\raisebox{0pt}[5mm][2mm]{}%
%EndExpansion
\emph{Abstract:}} \\ 
\multicolumn{7}{||l||}{%
%TCIMACRO{\TeXButton{Box}{\raisebox{0pt}[1cm][4cm]{}}}%
%BeginExpansion
\raisebox{0pt}[1cm][4cm]{}%
%EndExpansion
%TCIMACRO{%
%\TeXButton{AbstractBox}{\parbox{14.5cm}{The primary object of this report is to describe how a digital, causal,
%IIR, low-pass, high-pass, band-pass and band-stop filter may be designed from a set of
%specifications. As filter types, Butterworth and Chebyshev-1 are chosen.
%The design process is accomplished in three steps. First the parameters for an analog
%profile are deduced, then an analog profile with unity as the cut-off frequency
%calculated. Finally the decimal coefficients of a digital filter are obtained by mapping
%back this profile. In steps 1 and 3, the famous Bilinear Transformation is applied,
%and the method called Direct Calculations.
%
%Often, the efficiency of a microprocessor may be increased by restricting calculations
%on integer numbers. To obtain integer coefficients the above results are scaled
%properly and both the decimal and integer coefficients saved into text files. The integer
%coefficients are then used to construct a truncated impulse-response and the metod
%called Backward Calculations.
%
%The magnitude response of both methods may then be obtained and compared.
%We test the Direct against the Backward method for all filter functions and both
%filter types, i.e., low-pass, high-pass, band-pass and bans-stop as well as for the
%Butterworth and Chebyshev-1 types.
%
%The report then concludes with some comments on the results.}}}%
%BeginExpansion
\parbox{14.5cm}{The primary object of this report is to describe how a digital, causal,
IIR, low-pass, high-pass, band-pass and band-stop filter may be designed from a set of
specifications. As filter types, Butterworth and Chebyshev-1 are chosen.
The design process is accomplished in three steps. First the parameters for an analog
profile are deduced, then an analog profile with unity as the cut-off frequency
calculated. Finally the decimal coefficients of a digital filter are obtained by mapping
back this profile. In steps 1 and 3, the famous Bilinear Transformation is applied,
and the method called Direct Calculations.

Often, the efficiency of a microprocessor may be increased by restricting calculations
on integer numbers. To obtain integer coefficients the above results are scaled
properly and both the decimal and integer coefficients saved into text files. The integer
coefficients are then used to construct a truncated impulse-response and the metod
called Backward Calculations.

The magnitude response of both methods may then be obtained and compared.
We test the Direct against the Backward method for all filter functions and both
filter types, i.e., low-pass, high-pass, band-pass and bans-stop as well as for the
Butterworth and Chebyshev-1 types.

The report then concludes with some comments on the results.}%
%EndExpansion
} \\ \hline\hline
\end{tabular}

%TCIMACRO{\TeXButton{clearpage}{\clearpage}}%
%BeginExpansion
\clearpage%
%EndExpansion

\begin{tabular}{||p{0.8cm}|p{13.5cm}||}
\hline\hline
\multicolumn{2}{||c||}{%
%TCIMACRO{\TeXButton{Box}{\raisebox{0pt}[5mm][2mm]{}}}%
%BeginExpansion
\raisebox{0pt}[5mm][2mm]{}%
%EndExpansion
REVISION\ RECORD} \\ \hline
\multicolumn{1}{||c|}{%
%TCIMACRO{\TeXButton{Box}{\raisebox{0pt}[5mm][2mm]{}}}%
%BeginExpansion
\raisebox{0pt}[5mm][2mm]{}%
%EndExpansion
\emph{Rev}} & \multicolumn{1}{|c||}{\emph{Description}} \\ \hline
\multicolumn{1}{||c|}{%
%TCIMACRO{\TeXButton{Box}{\raisebox{0pt}[10mm][7mm]{}}}%
%BeginExpansion
\raisebox{0pt}[10mm][7mm]{}%
%EndExpansion
\emph{1}} & \multicolumn{1}{|c||}{%
%TCIMACRO{\TeXButton{parbox}{\parbox{13cm}{Original issue}}}%
%BeginExpansion
\parbox{13cm}{Original issue}%
%EndExpansion
} \\ \hline
\multicolumn{1}{||c|}{%
%TCIMACRO{\TeXButton{Box}{\raisebox{0pt}[10mm][7mm]{}}}%
%BeginExpansion
\raisebox{0pt}[10mm][7mm]{}%
%EndExpansion
\emph{2}} & \multicolumn{1}{|c||}{%
%TCIMACRO{\TeXButton{parbox}{\parbox{13cm}{}}}%
%BeginExpansion
\parbox{13cm}{}%
%EndExpansion
} \\ \hline
\multicolumn{1}{||c|}{%
%TCIMACRO{\TeXButton{Box}{\raisebox{0pt}[10mm][7mm]{}}}%
%BeginExpansion
\raisebox{0pt}[10mm][7mm]{}%
%EndExpansion
\emph{3}} & \multicolumn{1}{|c||}{%
%TCIMACRO{\TeXButton{parbox}{\parbox{13cm}{}}}%
%BeginExpansion
\parbox{13cm}{}%
%EndExpansion
} \\ \hline
\multicolumn{1}{||c|}{%
%TCIMACRO{\TeXButton{Box}{\raisebox{0pt}[10mm][7mm]{}}}%
%BeginExpansion
\raisebox{0pt}[10mm][7mm]{}%
%EndExpansion
\emph{4}} & \multicolumn{1}{|c||}{%
%TCIMACRO{\TeXButton{parbox}{\parbox{13cm}{}}}%
%BeginExpansion
\parbox{13cm}{}%
%EndExpansion
} \\ \hline
\multicolumn{1}{||c|}{%
%TCIMACRO{\TeXButton{Box}{\raisebox{0pt}[10mm][7mm]{}}}%
%BeginExpansion
\raisebox{0pt}[10mm][7mm]{}%
%EndExpansion
\emph{5}} & \multicolumn{1}{|c||}{%
%TCIMACRO{\TeXButton{parbox}{\parbox{13cm}{}}}%
%BeginExpansion
\parbox{13cm}{}%
%EndExpansion
} \\ \hline
\multicolumn{1}{||c|}{%
%TCIMACRO{\TeXButton{Box}{\raisebox{0pt}[10mm][7mm]{}}}%
%BeginExpansion
\raisebox{0pt}[10mm][7mm]{}%
%EndExpansion
\emph{6}} & \multicolumn{1}{|c||}{%
%TCIMACRO{\TeXButton{parbox}{\parbox{13cm}{}}}%
%BeginExpansion
\parbox{13cm}{}%
%EndExpansion
} \\ \hline
\multicolumn{1}{||c|}{%
%TCIMACRO{\TeXButton{Box}{\raisebox{0pt}[10mm][7mm]{}}}%
%BeginExpansion
\raisebox{0pt}[10mm][7mm]{}%
%EndExpansion
\emph{7}} & \multicolumn{1}{|c||}{%
%TCIMACRO{\TeXButton{parbox}{\parbox{13cm}{}}}%
%BeginExpansion
\parbox{13cm}{}%
%EndExpansion
} \\ \hline
\multicolumn{1}{||c|}{%
%TCIMACRO{\TeXButton{Box}{\raisebox{0pt}[10mm][7mm]{}}}%
%BeginExpansion
\raisebox{0pt}[10mm][7mm]{}%
%EndExpansion
\emph{8}} & \multicolumn{1}{|c||}{%
%TCIMACRO{\TeXButton{parbox}{\parbox{13cm}{}}}%
%BeginExpansion
\parbox{13cm}{}%
%EndExpansion
} \\ \hline
\multicolumn{1}{||c|}{%
%TCIMACRO{\TeXButton{Box}{\raisebox{0pt}[10mm][7mm]{}}}%
%BeginExpansion
\raisebox{0pt}[10mm][7mm]{}%
%EndExpansion
\emph{9}} & \multicolumn{1}{|c||}{%
%TCIMACRO{\TeXButton{parbox}{\parbox{13cm}{}}}%
%BeginExpansion
\parbox{13cm}{}%
%EndExpansion
} \\ \hline
\multicolumn{1}{||c|}{%
%TCIMACRO{\TeXButton{Box}{\raisebox{0pt}[10mm][7mm]{}}}%
%BeginExpansion
\raisebox{0pt}[10mm][7mm]{}%
%EndExpansion
\emph{10}} & \multicolumn{1}{|c||}{%
%TCIMACRO{\TeXButton{parbox}{\parbox{13cm}{}}}%
%BeginExpansion
\parbox{13cm}{}%
%EndExpansion
} \\ \hline
\multicolumn{1}{||c|}{%
%TCIMACRO{\TeXButton{Box}{\raisebox{0pt}[10mm][7mm]{}}}%
%BeginExpansion
\raisebox{0pt}[10mm][7mm]{}%
%EndExpansion
\emph{11}} & \multicolumn{1}{|c||}{%
%TCIMACRO{\TeXButton{parbox}{\parbox{13cm}{}}}%
%BeginExpansion
\parbox{13cm}{}%
%EndExpansion
} \\ \hline
&  \\ \hline\hline
\end{tabular}
%TCIMACRO{\TeXButton{clearpage}{\clearpage}}%
%BeginExpansion
\clearpage%
%EndExpansion

\QTP{contents}
TABLE\ OF\ CONTENTS

%TCIMACRO{\TeXButton{TableOfContents}{\TableOfContents} }%
%BeginExpansion
\TableOfContents
%EndExpansion
%TCIMACRO{\TeXButton{clearpage}{\clearpage}}%
%BeginExpansion
\clearpage%
%EndExpansion

\section{INTRODUCTION\label{Introduction}\protect\ref{Introduction}}

In this article we describe in detail the mathematical background of digital
IIR filter design. As filter types Butterworth and Chebyshev-1 analog
filters are chosen. To make the process more specific we design first a
digital low-pass filter and show that a high-pass, band-pass and band-stop
filter may be deduced by a frequency transformation. Then we normalize the
filter coefficients such that the magnitude response is in the interval $%
[0,\,1]$. To increase the efficiency of some microprocessors the decimal
integers are scaled properly to integers. By using these coefficients in a
method called Backward Calculations, we may then compare the magnitude
responses.

Denote the input, output and impulse-response of a linear, time-invariant
system, abbreviated LTI, by $x\left[ n\right] $, $y\left[ n\right] $ and $h%
\left[ n\right] $. Then the convolution relates these variables as 
\begin{gather}
y=h\ast x=x\ast h, \\
y\left[ n\right] =\sum_{k=-\infty }^{\infty }h\left[ k\right] x\left[ n-k%
\right] =\sum_{k=-\infty }^{\infty }x\left[ k\right] h\left[ n-k\right] .
\end{gather}

Denote by $x_{a}(t)$ a continuous function and construct a discrete function 
$x[n]$ by sampling $x_{a}(t)$ with sampling period $T$. Thus 
\begin{equation}
x\left[ n\right] =x_{a}\left( n\,T\right) .
\end{equation}
Denote by $\mathcal{L}$, $\mathcal{Z}$ and $\mathcal{F}$ the operations of 
\textit{Laplace, }$Z$ and \textit{Fourier-}transformations given by 
\begin{gather}
X_{a}\left( s\right) =\mathcal{L}\left\{ x_{a}\left( t\right) \right\}
=\int_{-\infty }^{\infty }x_{a}\left( t\right) \,e^{-s\,t}\,dt,
\label{Laplace_Transform} \\
X\left( z\right) =\mathcal{Z}\left\{ \,x\left[ n\right] \,\right\}
=\sum_{k=-\infty }^{\infty }x\left[ k\right] z^{-k},  \label{Z_Transform} \\
X\left( e^{j\,\omega }\right) =\mathcal{F}\left\{ \,x\left[ n\right]
\,\right\} =\sum_{k=-\infty }^{\infty }x\left[ k\right] e^{-j\,\omega \,k}.
\label{Fourier_Transform}
\end{gather}
Comparing Eqs. (\ref{Z_Transform}) and (\ref{Fourier_Transform}) we see that
the $\mathcal{F}$-transform of a sequence is obtained from its $\mathcal{Z}$%
-transform by replacing the complex variable $z$ with $e^{j\,\omega }$. In a
linear, time-invariant system the input and output variables are related by 
\begin{gather}
Y\left( z\right) =H\left( z\right) X\left( z\right) , \\
Y\left( e^{j\,\omega }\right) =H\left( e^{j\,\omega }\right) X\left(
e^{j\,\omega }\right) ,
\end{gather}
where $H\left( z\right) $ and $Y\left( z\right) $ are $\mathcal{Z}$%
-transforms of $h\left[ n\right] $, while $H\left( e^{jw}\right) $ and $%
Y\left( e^{jw}\right) $ are $\mathcal{F}$-transforms of $h\left[ n\right] $
and $y\left[ n\right] $, respectively. From Eqs. (\ref{Laplace_Transform}), (%
\ref{Z_Transform}) and (\ref{Fourier_Transform}) we see that for a
continuous or discrete, real function in the time domain, the complex
conjugate transform is obtained by replacing the independent variable with
its complex conjugate counterpart, 
\begin{eqnarray}
x_{a}\left( t\right) \;\func{real} &\leftrightarrow &\left| X_{a}\left(
\sigma +j\,\omega \right) \right| ^{2}=X_{a}\left( \sigma +j\,\omega \right)
X_{a}\left( \sigma -j\,\omega \right) , \\
x\left[ n\right] \;\func{real} &\leftrightarrow &\left| X\left(
r\,e^{j\,\omega }\right) \right| ^{2}=X\left( r\,e^{j\,\omega }\right)
X\left( r\,e^{-j\,\omega }\right) .
\end{eqnarray}

An Infinite Impulse-Response, abbreviated IIR, causal filter of order $N$
may be realized as an $N$th order difference Eq. of the form 
\begin{equation}
y\left[ n\right] =\sum_{k=0}^{N}a_{k}x\left[ n-k\right] +\sum_{k=1}^{N}b_{k}y%
\left[ n-k\right] .  \label{LTI_Causal_Filter}
\end{equation}

Thus, to completely determine an $N$th order IIR filter, we have to
calculate two vectors $\vec{a}$ and $\vec{b}$ of length $N+1$, with $b_{0}=0$%
. Applying the $\mathcal{Z}$ and $\mathcal{F}$-transformations to Eq. (\ref%
{LTI_Causal_Filter}), its system function and frequency-response can be
obtained as 
\begin{gather}
H\left( z\right) =\frac{Y\left( z\right) }{X\left( z\right) }=\frac{%
\sum_{k=0}^{N}a_{k}\,z^{-k}}{1-\sum_{k=1}^{N}b_{k}\,z^{-k}},
\label{System_Function} \\
H\left( e^{j\,\omega }\right) =\frac{Y\left( e^{j\,\omega }\right) }{X\left(
e^{j\,\omega }\right) }=\frac{\sum_{k=0}^{N}a_{k}\,e^{-j\,\omega \,k}}{%
1-\sum_{k=1}^{N}b_{k}\,e^{-j\,\omega \,k}}.  \label{Freq_Resp}
\end{gather}
It is customary to use $dB$-scale when attenuation of a frequency-response
variable is concerned, i.e. 
\begin{equation}
dB\text{-amplitude of }H\left( e^{j\,\omega }\right) =10\,\log _{10}\left|
H\left( e^{j\,\omega }\right) \right| ^{2}=20\,\log _{10}\left| H\left(
e^{j\,\omega }\right) \right| .
\end{equation}
\FRAME{ftbphFU}{12.2769cm}{9.228cm}{0pt}{\Qcb{Specification of a digital
low-pass filter.}}{\Qlb{Gen_LowPass}}{general_lowpass.eps}{\special{language
"Scientific Word";type "GRAPHIC";maintain-aspect-ratio TRUE;display
"ICON";valid_file "F";width 12.2769cm;height 9.228cm;depth
0pt;original-width 8.003in;original-height 6.0105in;cropleft "0";croptop
"0.9996";cropright "1.0007";cropbottom "0";filename
'General_LowPass.eps';file-properties "XNPEU";}}

\begin{center}
\FRAME{ftbphFU}{12.2681cm}{9.2324cm}{0pt}{\Qcb{Specification of a digital
high-pass filter.}}{\Qlb{Gen_HighPass}}{general_highpass.eps}{\special%
{language "Scientific Word";type "GRAPHIC";maintain-aspect-ratio
TRUE;display "ICON";valid_file "F";width 12.2681cm;height 9.2324cm;depth
0pt;original-width 8.003in;original-height 6.0105in;cropleft "0";croptop
"1";cropright "1";cropbottom "0";filename
'General_HighPass.eps';file-properties "XNPEU";}}

\FRAME{ftbphFU}{12.2681cm}{9.2302cm}{0pt}{\Qcb{Specification of a digital
band-pass filter.}}{\Qlb{Gen_BandPass}}{general_bandpass.eps}{\special%
{language "Scientific Word";type "GRAPHIC";maintain-aspect-ratio
TRUE;display "ICON";valid_file "F";width 12.2681cm;height 9.2302cm;depth
0pt;original-width 8.003in;original-height 6.0105in;cropleft "0";croptop
"1";cropright "1";cropbottom "0";filename
'General_BandPass.eps';file-properties "XNPEU";}}

\FRAME{ftbphFU}{12.2681cm}{9.2302cm}{0pt}{\Qcb{Specification of a digital
band-stop filter.}}{\Qlb{Gen_BandStop}}{general_bandstop.eps}{\special%
{language "Scientific Word";type "GRAPHIC";maintain-aspect-ratio
TRUE;display "ICON";valid_file "F";width 12.2681cm;height 9.2302cm;depth
0pt;original-width 8.003in;original-height 6.0105in;cropleft "0";croptop
"1";cropright "1";cropbottom "0";filename
'General_BandStop.eps';file-properties "XNPEU";}}A digital low-pass filter,
may then be specified by demanding that the frequencies lower than $\omega
_{1}$ must be attenuated by a factor $0\leq A\leq A_{1}$, while frequencies
higher than $\omega _{2}$ is attenuated by $A_{2}\leq A$. Thus, the design
specifications for a normalized, digital, low-pass filter are of the form 
\begin{gather}
0<\omega _{1}<\omega _{2}<\pi ,  \label{Spec_1_low} \\
0<A_{1}<A_{2},  \label{Spec_2_low} \\
-A_{1}\leq 20\,\log _{10}\left| H(e^{j\,\omega })\right| \leq 0,\text{ \ \ \
\ }0\leq \omega \leq \omega _{1},  \label{Spec_3_low} \\
-\infty <20\,\log _{10}\left| H(e^{j\,\omega })\right| \leq -A_{2},\text{ \
\ \ \ }\omega _{2}\leq \omega \leq \pi .  \label{Spec_4_low}
\end{gather}
\end{center}

A high-pass filter may be specified by requiring that frequencies up to $%
\omega _{1}$ is attenuated by $A_{1}\leq A$, and frequencies higher than $%
\omega _{2}$ is attenuated by $0\leq A\leq A_{2}$. Thus, a normalized
high-pass filter may be specified by design specifications of the form 
\begin{gather}
0<\omega _{1}<\omega _{2}<\pi ,  \label{Spec_1_high} \\
0<A_{2}<A_{1},  \label{Spec_2_high} \\
-\infty <20\,\log _{10}\left| H(e^{j\,\omega })\right| \leq -A_{1},\text{ \
\ \ \ }0\leq \omega \leq \omega _{1}.  \label{Spec_3_high} \\
-A_{2}\leq 20\,\log _{10}\left| H(e^{j\,\omega })\right| \leq 0,\text{ \ \ \
\ }\omega _{2}\leq \omega \leq \pi .  \label{Spec_4_high}
\end{gather}
A band-pass filter may be specified by requiring that frequencies up to $%
\omega _{1}$ is attenuated by $A_{1}\leq A$, frequencies between $\omega
_{2} $ and $\omega _{3}$ by $A\leq \min \left( A_{2},\,A_{3}\right) $ and
frequencies higher than $\omega _{4}$ by $A_{4}\leq A.$ Thus, a normalized
band-pass filter may be specified by design specifications of the form 
\begin{gather}
0<\omega _{1}<\omega _{2}<\omega _{3}<\omega _{4}<\pi ,
\label{Spec_1_BandPass} \\
0<A_{2}<A_{1},  \label{Spec_2_BandPass} \\
0<A_{3}<A_{4},  \label{Spec_3_BandPass} \\
-\infty <20\,\log _{10}\left| H(e^{j\,\omega })\right| \leq -A_{1},\text{ \
\ \ \ }0\leq \omega \leq \omega _{1},  \label{Spec_4_BandPass} \\
-\min \left( A_{2},\,A_{3}\right) \leq 20\,\log _{10}\left| H(e^{j\,\omega
})\right| \leq 0,\text{ \ \ \ \ }\omega _{2}\leq \omega \leq \omega _{3},
\label{Spec_5_BandPass} \\
-\infty <20\,\log _{10}\left| H(e^{j\,\omega })\right| \leq -A_{4},\text{ \
\ \ \ }\omega _{4}\leq \omega \leq \pi .  \label{Spec_6_BandPass}
\end{gather}
Finally, a band-stop filter may be specified by requiring that frequencies
up to $\omega _{1}$ is attenuated by $A\leq A_{1}$, frequencies between $%
\omega _{2}$ and $\omega _{3}$ by $\max \left( A_{2},\,A_{3}\right) \leq A$
and frequencies higher than $\omega _{4}$ by $A\leq A_{4}.$ Thus, a
normalized band-stop filter may be specified by design specifications of the
form 
\begin{gather}
0<\omega _{1}<\omega _{2}<\omega _{3}<\omega _{4}<\pi ,
\label{Spec_1_BandStop} \\
0<A_{1}<A_{2},  \label{Spec_2_BandStop} \\
0<A_{4}<A_{3},  \label{Spec_3_BandStop} \\
-A_{1}\leq 20\,\log _{10}\left| H(e^{j\,\omega })\right| \leq 0,\text{ \ \ \
\ }0\leq \omega \leq \omega _{1},  \label{Spec_4_BandStop} \\
-\infty \leq 20\,\log _{10}\left| H(e^{j\,\omega })\right| \leq -\max \left(
A_{2},\,A_{3}\right) ,\text{ \ \ \ \ }\omega _{2}\leq \omega \leq \omega
_{3},  \label{Spec_5_BandStop} \\
-A_{4}\leq 20\,\log _{10}\left| H(e^{j\,\omega })\right| \leq 0,\text{ \ \ \
\ }\omega _{4}\leq \omega \leq \pi .  \label{Spec_6_BandStop}
\end{gather}
Figs. \ref{Gen_LowPass} , \ref{Gen_HighPass}, \ref{Gen_BandPass} and \ref%
{Gen_BandStop} illustrate design specifications for low-pass, high-pass,
band-pass and a band-stop digital filters. Note that we have normalized the
frequency variable $f=\frac{\omega }{2\,\pi }$, with respect to the sampling
frequency. Furthermore, we have the following relations between values of $%
A_{i}$, $i=1,\ldots ,\,4$ and $\varepsilon _{i}$, $i=1,\,2,\,3$:

\begin{itemize}
\item For a low-pass filter 
\begin{eqnarray}
A_{1} &=&-20\,\log _{10}\left( 1-\varepsilon _{1}\right) , \\
A_{2} &=&-20\,\log _{10}\left( \varepsilon _{2}\right) .
\end{eqnarray}

\item For a high-pass filter 
\begin{eqnarray}
A_{1} &=&-20\,\log _{10}\left( \varepsilon _{2}\right) , \\
A_{2} &=&-20\,\log _{10}\left( 1-\varepsilon _{1}\right) .
\end{eqnarray}

\item For a band-pass filter 
\begin{eqnarray}
A_{1} &=&-20\,\log _{10}\left( \varepsilon _{2}\right) , \\
A_{2},\,A_{3} &\leq &-20\,\log _{10}\left( 1-\varepsilon _{1}\right) , \\
A_{4} &=&-20\,\log _{10}\left( \varepsilon _{3}\right) .
\end{eqnarray}

\item For a band-stop filter 
\begin{eqnarray}
A_{1},\,A_{4} &\leq &-20\,\log _{10}\left( 1-\varepsilon _{1}\right) , \\
A_{2} &=&-20\,\log _{10}\left( \varepsilon _{2}\right) , \\
A_{3} &=&-20\,\log _{10}\left( \varepsilon _{3}\right) .
\end{eqnarray}
\end{itemize}

\begin{remark}
By comparing the low-pass with high-pass specifications, Eqs. (\ref%
{Spec_1_low}) through (\ref{Spec_4_low}) with (\ref{Spec_1_high}) through (%
\ref{Spec_4_high}) or Fig \ref{Gen_LowPass} with \ref{Gen_HighPass}, we
observe that a low-pass filter may be derived from a high-pass filter by
changing the frequency variable from $\omega _{hp}$ to $\omega _{lp}=\pi
-\omega _{hp}.$ Specifically, 
\begin{eqnarray}
\omega _{1,\,lp} &=&\pi -\omega _{2,\,hp}, \\
\omega _{2,\,lp} &=&\pi -\omega _{1,\,hp}, \\
A_{2,\,lp} &=&A_{1,\,hp}, \\
A_{1,\,lp} &=&A_{2,\,hp}.
\end{eqnarray}
\end{remark}

\begin{remark}
The band-pass and band-stop specifications in Eqs. (\ref{Spec_1_BandPass})
through (\ref{Spec_6_BandStop}) do not need to be symmetrical. This is
clearly seen for example in the specification of the form 
\begin{gather}
\omega _{2}-\omega _{1}=\omega _{4}-\omega _{3}, \\
\begin{array}{ll}
A_{1}\neq A_{4}, & \text{for a band-pass filter,} \\ 
A_{2}\neq A_{3}. & \text{for a band-stop filter,}%
\end{array}%
\end{gather}
or from Figs. \ref{Gen_BandPass} and \ref{Gen_BandStop}.
\end{remark}

\begin{remark}
By comparing the low-pass and high-pass specifications with band-pass
specifications, Eqs. (\ref{Spec_1_low}) through (\ref{Spec_4_high}) with (%
\ref{Spec_1_BandPass}) through (\ref{Spec_4_BandPass}) or Fig \ref%
{Gen_LowPass} and \ref{Gen_HighPass} with \ref{Gen_BandPass}, we observe
that a band-pass filter may be derived by cascading a low-pass and a
high-pass filter. As we see later in this report we will not apply this
approach to implement a band-pass filter.
\end{remark}

\begin{remark}
A band-pass filter specification may be thought as a combination of a
high-pass and a low-pass specification. Specifically 
\begin{equation}
\begin{array}{ll}
\omega _{i,\,bp}=\omega _{i,\,hp}, & i=1,2, \\ 
\omega _{i,\,bp}=\omega _{i-2,\,lp}, & i=3,4, \\ 
A_{i,\,bp}=A_{i,\,hp}, & i=1,2, \\ 
A_{i,\,bp}=A_{i-2,\,lp}, & i=3,4.%
\end{array}%
\end{equation}
\end{remark}

\begin{remark}
Compare now low-pass and high-pass specifications with band-stop
specifications, Eqs. (\ref{Spec_1_low}) through (\ref{Spec_4_high}) with (%
\ref{Spec_1_BandStop}) through (\ref{Spec_4_BandStop}) or Fig \ref%
{Gen_LowPass} and \ref{Gen_HighPass} with \ref{Gen_BandStop}. Observe that a
band-stop filter specification is composed of first a low-pass then a
high-pass filter specification. Specifically 
\begin{equation}
\begin{array}{ll}
\omega _{i,\,bp}=\omega _{i,\,lp}, & i=1,2, \\ 
\omega _{i,\,bp}=\omega _{i-2,\,hp}, & i=3,4, \\ 
A_{i,\,bp}=A_{i,\,lp}, & i=1,2, \\ 
A_{i,\,bp}=A_{i-2,\,hp}, & i=3,4.%
\end{array}%
\end{equation}
\end{remark}

We take advantage of \ these properties later to determine the high-pass,
band-pass and band-stop filter parameters.

For a detailed discussion of the above matters, consult Chapters 2, 3, 4 5,
7 and 8 of \cite{Oppenheim}.

\section{ANALOG FILTER TYPES\protect\nolinebreak}

\subsection{\protect\nolinebreak BUTTERWORTH}

The squared response of an analog, $N$th order, low-pass, Butterworth filter
is defined by 
\begin{equation}
\left| H_{a}\left( j\,\Omega \right) \right| ^{2}=H_{a}\left( s\right)
H_{a}\left( -s\right) =\frac{1}{1+\left( -s^{2}\right) ^{N}}.
\label{Butterworth}
\end{equation}
By analytic continuation Eq. (\ref{Butterworth}) can be extended to the
complex $s$-domain giving 
\begin{equation}
\left| H_{a}\left( j\,\Omega \right) \right| ^{2}=H_{a}\left( j\,\Omega
\right) H_{a}\left( -j\,\Omega \right) =\frac{1}{1+\left( \Omega /\Omega
_{c}\right) ^{2\,N}}.  \label{Butterworth2}
\end{equation}
Poles of the squared magnitude response may thus be found as 
\begin{equation}
s_{k}=\Omega _{c}\,e^{j\,\pi \left[ \frac{1}{2}+\frac{2\,k+1}{2\,N}\right] },%
\text{ \ \ \ \ }k=0,\,1,\,\ldots ,\,2N-1.
\end{equation}
The poles lie equally spaced on a circle of radius $\Omega _{c}$ centered
about $s=0$. Thus 
\begin{gather}
s_{k}=\sigma _{k}+j\,\Omega _{k},\text{ \ \ \ \ }k=0,\,1,\,2,\,\ldots
,\,2N-1,  \label{Butter_Poles} \\
\sigma _{k}=-\Omega _{c}\sin \left[ \frac{\left( 2k+1\right) \pi }{2N}\right]
,  \label{Butter_Real_Part} \\
\Omega _{k}=-\Omega _{c}\cos \left[ \frac{\left( 2k+1\right) \pi }{2N}\right]
.  \label{Butter_Imag_Part}
\end{gather}
A Butterworth low-pass filter is maximally flat at $\Omega =0$; that is the
first $2N-1$ derivatives of $\left| H_{a}\left( j\,\Omega \right) \right|
^{2}$ at $\Omega =0$ are equal to zero. Note that $\left| H_{a}(j\,\Omega
)\right| _{\Omega =\Omega _{c}}^{2}=\frac{1}{2}$, regardless of order $N$,
so the magnitude is down by $3\,dB$ at $\Omega =\Omega _{c}$. Also the
magnitude is a monotonically decreasing function of $\left| \Omega \right| $%
. It can be seen from Figs. \ref{Low_Butter_Plot} that as $N$ increases the
following changes occur in the magnitude characteristics:

\begin{enumerate}
\item The transition from the pass-band to the stop-band becomes more rapid.

\item The magnitude $\left| H_{a}\left( j\,\Omega \right) \right| \approx 1$
for more of the pass-band.

\item The magnitude $\left| H_{a}\left( j\,\Omega \right) \right| \approx 0$
for more of the stop-band.
\end{enumerate}

\subsection{CHEBYSHEV}

Chebyshev filters relax the constraint of monotonicity over either the
pass-band or the stop-band. For a given filter order $N$, the Chebyshev
filter gives

\begin{itemize}
\item minimum equiripple deviation of the magnitude characteristic over one
prescribed band of frequencies and

\item monotonic behavior over the remaining band of frequencies.
\end{itemize}

Depending on the design choice, the minimum equiripple error property can be
associated with either the pass-band (Chebyshev-1 filter) or the stop-band
(Chebyshev-2 filter). For a given set of filter specifications, this
concession of allowing ripple in a band of frequencies leads to lower-order
than the Butterworth design choice. Chebyshev-2 filters are rarely used and
therefore not included here.

\subsubsection{CHEBYSHEV-1 Filters}

An analog, $N$th order, low-pass, Chebyshev-1 filter exhibits equiripple
error in the pass-band and a monotonically decreasing response in the
stop-band. It is specified by the squared-magnitude function 
\begin{equation}
\left| H_{a}\left( j\,\Omega \right) \right| ^{2}=\frac{1}{1+\epsilon
^{2}T_{N}^{2}\left( \Omega /\Omega _{1}\right) },  \label{Cheb1_Low_Filter}
\end{equation}
where $T_{n}\left( x\right) $, $n=0,\,1,\,2,\,\ldots $ is the $n$th order
Chebyshev polynomial given by 
\begin{equation}
T_{n}\left( x\right) =\left\{ 
\begin{array}{c}
\cos \left( n\,\cos ^{-1}x\right) ,\text{ \ \ \ \ }\left| x\right| \leq 1,
\\ 
\cosh \left( n\,\cosh ^{-1}x\right) ,\text{ \ \ \ \ }\left| x\right| \geq 1.%
\end{array}
\right.
\end{equation}
This polynomial can be derived from the recurrence relation 
\begin{eqnarray}
T_{0}\left( x\right) &=&1, \\
T_{1}\left( x\right) &=&x, \\
T_{n+1}\left( x\right) &=&2x\,T_{n}\left( x\right) -T_{n-1}\left( x\right) ,%
\text{ \ \ \ \ }n=1,2,\,\ldots .
\end{eqnarray}
From the Chebyshev polynomials definition we see that 
\begin{eqnarray}
\left| T_{n}\left( x\right) \right| &\leq &1,\text{ \ \ \ \ }\left| x\right|
\leq 1,  \label{ChebProp1} \\
\left| T_{n}\left( x\right) \right| &\geq &1,\text{ \ \ \ \ }\left| x\right|
\geq 1.  \label{ChebProp2}
\end{eqnarray}
Furthermore, the $N$th order Chebyshev polynomial gives the smallest maximum
value of all $N$th order polynomials. In particular, observe from Eq. (\ref%
{Cheb1_Low_Filter}) that $\frac{1}{1+\epsilon ^{2}}$ $\leq \left|
H_{a}\left( j\,\Omega \right) \right| ^{2}\leq 1$ within the pass-band,
where it will have a total of $N$ local maxima and local minima. Observe
from Eq. (\ref{Cheb1_Low_Filter}) that Chebyshev-1 low-pass filters with odd 
$N$ has $\left| H\left( j\,\Omega \right) \right| _{\Omega =0}^{2}=1$, while
Chebyshev-1 low-pass filters with even $N$ has $\left| H\left( j\,\Omega
\right) \right| _{\Omega =0}^{2}=\frac{1}{1+\epsilon ^{2}}$. As we see
later, this value is used to normalize the digital low-pass filters. All
Chebyshev-1 low-pass filters has $\left| H\left( j\,\Omega \right) \right|
_{_{\Omega =\Omega _{p}}}^{2}=\frac{1}{1+\epsilon ^{2}}$. The location of
the poles are obtained from Eq. (\ref{Cheb1_Low_Filter}) as 
\begin{gather}
s_{k}=\sigma _{k}+j\,\Omega _{k},\text{ \ \ \ \ }k=0,\,1,\,2,\,\ldots
,\,2N-1,  \label{Cheb_Poles} \\
\sigma _{k}=-\mu \,\Omega _{1}\sin \left[ \frac{\left( 2k+1\right) \pi }{2N}%
\right] ,  \label{Cheb_Real_Part} \\
\Omega _{k}=\nu \,\Omega _{1}\cos \left[ \frac{\left( 2k+1\right) \pi }{2N}%
\right] ,  \label{Cheb_Imag_Part} \\
\mu =\frac{1}{2}\left( \gamma -\gamma ^{-1}\right) ,  \label{Minor_Axis} \\
\nu =\frac{1}{2}\left( \gamma +\gamma ^{-1}\right) ,  \label{Major_Axis} \\
\gamma =\left( \frac{1+\sqrt{1+\epsilon ^{2}}}{\epsilon }\right) ^{1/N}.
\label{Gamma}
\end{gather}
As we see from Eqs. (\ref{Cheb_Real_Part}) through (\ref{Gamma}) these poles
are equally spaced about the origin on an ellipse with minor-axis length $%
\mu \,\Omega _{1}$ and major-axis length $\nu \,\Omega _{1}$ in the $s$%
-plane. Thus 
\begin{equation}
\frac{\sigma _{k}^{2}}{\left( \mu \,\Omega _{1}\right) ^{2}}+\frac{\Omega
_{k}^{2}}{\left( \nu \,\Omega _{1}\right) ^{2}}=1.
\end{equation}
Fig.\ref{Low_Butter_Plot} and \ref{Low_Cheb_Plot} show squared responses of
analog low-pass Butterworth and Chebyshev-1 filters.

\begin{center}
\FRAME{ftbphFU}{11.2489cm}{8.468cm}{0pt}{\Qcb{Squared magnitude responses
for analog low-pass Butterworth filters with $\Omega _{c}=2$.}}{\Qlb{%
Low_Butter_Plot}}{low_butter.eps}{\special{language "Scientific Word";type
"GRAPHIC";maintain-aspect-ratio TRUE;display "ICON";valid_file "F";width
11.2489cm;height 8.468cm;depth 0pt;original-width 8.003in;original-height
6.0105in;cropleft "0";croptop "1";cropright "1";cropbottom "0";filename
'Low_Butter.eps';file-properties "XNPEU";}}

\FRAME{ftbphFU}{11.2445cm}{8.468cm}{0pt}{\Qcb{Squared magnitude responses of
analog low-pass Chebyshev-1 filters with $\Omega _{1}=2$ and $\Omega _{2}=3$.%
}}{\Qlb{Low_Cheb_Plot}}{low_cheb.eps}{\special{language "Scientific
Word";type "GRAPHIC";maintain-aspect-ratio TRUE;display "ICON";valid_file
"F";width 11.2445cm;height 8.468cm;depth 0pt;original-width
8.003in;original-height 6.0105in;cropleft "0";croptop "1.0001";cropright
"0.9996";cropbottom "0";filename 'Low_Cheb.eps';file-properties "XNPEU";}}
\end{center}

For further details on these filter types \cite{Mitra}, sections 5-3-1 and
5-3-2 may be consulted.

\section{METHOD}

To design an IIR digital filter we begin by an analog profile and then map
it to a digital system function. This is because the theory of analog filter
design is well established and many closed-form design formulas exist. Given
the design specifications of a digital filter, the derivation of the digital
filter system function requires three steps:

\begin{itemize}
\item Map the desired digital filter specifications into those for an
equivalent analog filter.

\item Derive a corresponding analog filter transfer function for the analog
prototype.

\item Convert the transfer function of the analog prototype into an
equivalent digital filter system function.
\end{itemize}

\begin{remark}
Note that the above steps are not independent. Specifically, step 1
transforms the digital specifications to an analog counterpart, while step 3
converts the resulting analog prototype to a digital system function.\label%
{DigAnalDig}
\end{remark}

\subsection{BILINEAR TRANSFORMATION}

To transform an analog to a digital profile or vice versa, we use the most
popular method, namely the \textit{Bilinear Transformation}. It relates the
two prototypes through the following relations 
\begin{eqnarray}
H\left( z\right) &=&H_{a}\left( s\right) \left| _{s=\alpha \frac{1-z^{-1}}{%
1+z^{-1}}}\right. ,  \label{Bilinear_Trans} \\
H_{a}\left( s\right) &=&H\left( z\right) \left| _{z=\frac{1+s/\alpha }{%
1-s/\alpha }}\right. .  \label{Bilinear_Inverse_Trans}
\end{eqnarray}
The digital system function is independent of the value of $\alpha $, so it
may be chosen arbitrary. This is because step 3 is the inverse of step 1 in
the design procedure, see Remark (\ref{DigAnalDig}). The transformation in
Eq. (\ref{Bilinear_Trans}) maps the $s$-plane into the $z-$plane as follows:

\begin{enumerate}
\item Left-half $s$-plane\bigskip\ $\longrightarrow $ interior of $z-$plane
unit circle.

\item Right-half $s$-plane\bigskip\ $\longrightarrow $ exterior of $z-$plane
unit circle.

\item Imaginary axis, $s=j\,\Omega \longrightarrow $ $z-$plane unit circle, $%
z=e^{j\,\omega }$.
\end{enumerate}

Furthermore, the mapping in step 3 above is non-linear. Substituting $%
z=e^{j\,\omega }$ into Eq. (\ref{Bilinear_Trans}) results in 
\begin{equation}
H\left( e^{j\,\omega }\right) =H_{a}\left( j\,\alpha \,\tan \frac{\omega }{2}%
\right) ,
\end{equation}
which gives the relation between the frequency variables as 
\begin{eqnarray}
\Omega &=&\alpha \,\tan \left( \frac{\omega }{2}\right) ,
\label{Dig_2_Analog_Freq} \\
\omega &=&2\,\tan ^{-1}\left( \frac{\Omega }{\alpha }\right) .
\label{Anal_2_Dig_Freq}
\end{eqnarray}
For further details on the Bilinear transformation, refer to \cite{Smith},
pages 626-628 and \cite{Mitra}, section 5-4-1. To normalize the
transformation we have chosen $\alpha $ as 
\begin{equation}
\alpha =\frac{1}{\tan \left( \frac{1}{2}\right) },
\end{equation}
which maps $\Omega =1\longleftrightarrow \omega =1$. Eq. (\ref%
{Dig_2_Analog_Freq}) is commonly referred to as frequency prewarping. Thus
the entire Bilinear Transformation Method is as follows:

\begin{enumerate}
\item Apply Eq. (\ref{Dig_2_Analog_Freq}) to prewarp all critical digital
frequencies given in the filter specifications, i.e. $\omega _{1}$ and $%
\omega _{2}$.

\item Given the prewarped analog critical frequencies, derive appropriate
analog transfer function $H_{a}\left( s\right) $.

\item Apply the bilinear transformation, Eq. (\ref{Bilinear_Trans}), to
obtain the desired digital system function $H\left( z\right) $.
\end{enumerate}

\subsection{MAPPING DESIGN\ SPECIFICATIONS}

Remember that a low-pass and high-pass filter had 2 critical frequencies and
2 critical attenuations, while in the band-pass and band-stop case we need 4
critical frequencies and 4 critical attenuations. The analog critical
frequencies may be obtained from Eq. (\ref{Dig_2_Analog_Freq}).

\begin{itemize}
\item Design specifications for a low-pass filter, Eqs. (\ref{Spec_1_low})
through (\ref{Spec_4_low}) may be mapped as 
\begin{eqnarray}
\Omega _{i} &=&\alpha \,\tan \left( \omega _{i}/2\right) ,\;\;\;\;\;i=1,\,2,
\label{Analog_Spec_0_low} \\
0 &<&\Omega _{1}<\Omega _{2}<\infty ,  \label{Analog_Spec_1_low} \\
0 &<&A_{1}<A_{2}<\infty ,  \label{Analog_Spec_2_low} \\
-A_{1} &\leq &10\,\log _{10}\left| H(j\,\Omega )\right| ^{2}\leq 0,\text{ \
\ \ \ }0\leq \Omega \leq \Omega _{1},  \label{Analog_Spec_3_low} \\
-\infty &<&10\,\log _{10}\left| H(j\,\Omega )\right| ^{2}\leq -A_{2},\text{
\ \ \ \ }\Omega _{2}\leq \Omega <\infty .  \label{Analog_Spec_4_low}
\end{eqnarray}

\item An analog high-pass filter is specified by mapping Eqs. (\ref%
{Spec_1_high}) through (\ref{Spec_4_high}). Thus 
\begin{eqnarray}
\Omega _{i} &=&\alpha \,\tan \left( \omega _{i}/2\right) ,\;\;\;\;\;i=1,\,2,
\label{Analog_Spec_0_high} \\
0 &<&\Omega _{1}<\Omega _{2}<\infty ,  \label{Analog_Spec_1_high} \\
0 &<&A_{2}<A_{1}<\infty ,  \label{Analog_Spec_2_high} \\
-\infty &<&10\,\log _{10}\left| H(j\,\Omega )\right| ^{2}\leq -A_{1},\text{
\ \ \ \ }0\leq \Omega \leq \Omega _{1},  \label{Analog_Spec_3_high} \\
-A_{2} &\leq &10\,\log _{10}\left| H(j\,\Omega )\right| ^{2}\leq 0,\text{ \
\ \ \ }\Omega _{2}\leq \Omega <\infty .  \label{Analog_Spec_4_high}
\end{eqnarray}

\item An analog band-pass filter is specified by mapping Eqs. (\ref%
{Spec_1_BandPass}) through (\ref{Spec_5_BandPass}). Thus 
\begin{eqnarray}
\Omega _{i} &=&\alpha \,\tan \left( \omega _{i}/2\right)
,\;\;\;\;\;i=1,\,\ldots ,\,4,  \label{Analog_Spec_0_BandPass} \\
0 &<&\Omega _{1}<\Omega _{2}<\Omega _{3}<\Omega _{4}<\infty ,
\label{Analog_Spec_1_BandPass} \\
0 &<&A_{2}<A_{1}<\infty ,  \label{Analog_Spec_2_BandPass} \\
0 &<&A_{3}<A_{4}<\infty ,  \label{Analog_Spec_3_BandPass} \\
-\infty &<&10\,\log _{10}\left| H(j\,\Omega )\right| ^{2}\leq -A_{1},\text{
\ \ \ \ }0\leq \Omega \leq \Omega _{1},  \label{Analog_Spec_4_BandPass} \\
-\min \left( A_{2},\,A_{3}\right) &\leq &10\,\log _{10}\left| H(j\,\Omega
)\right| ^{2}\leq 0,\text{ \ \ \ \ }\Omega _{2}\leq \Omega \leq \Omega _{3},
\label{Analog_Spec_5_BandPass} \\
-\infty &<&10\,\log _{10}\left| H(j\,\Omega )\right| ^{2}\leq -A_{4},\text{
\ \ \ \ }\Omega _{4}\leq \Omega <\infty .  \label{Analog_Spec_6_BandPass}
\end{eqnarray}

\item An analog band-stop filter is specified by mapping Eqs. (\ref%
{Spec_1_BandStop}) through (\ref{Spec_5_BandStop}). Thus 
\begin{eqnarray}
\Omega _{i} &=&\alpha \,\tan \left( \omega _{i}/2\right)
,\;\;\;\;\;i=1,\,\ldots ,\,4,  \label{Analog_Spec_0_BandStop} \\
0 &<&\Omega _{1}<\Omega _{2}<\Omega _{3}<\Omega _{4}<\infty ,
\label{Analog_Spec_1_BandStop} \\
0 &<&A_{1}<A_{2}<\infty ,  \label{Analog_Spec_2_BandStop} \\
0 &<&A_{4}<A_{3}<\infty ,  \label{Analog_Spec_3_BandStop} \\
-A_{1} &<&10\,\log _{10}\left| H(j\,\Omega )\right| ^{2}\leq 0,\text{ \ \ \
\ }0\leq \Omega \leq \Omega _{1},  \label{Analog_Spec_4_BandStop} \\
-\infty &<&10\,\log _{10}\left| H(j\,\Omega )\right| ^{2}\leq -\max \left(
A_{2},\,A_{3}\right) ,\text{ \ \ \ \ }\Omega _{2}\leq \Omega \leq \Omega
_{3},  \label{Analog_Spec_5_BandStop} \\
-A_{4} &\leq &10\,\log _{10}\left| H(j\,\Omega )\right| ^{2}\leq 0,\text{ \
\ \ \ }\Omega _{4}\leq \Omega <\infty .  \label{Analog_Spec_6_BandStop}
\end{eqnarray}
\end{itemize}

\subsubsection{BUTTERWORTH}

From Eq. (\ref{Butterworth2}), we observe that to fully determine an analog,
low-pass, Butterworth filter we need to specify $\Omega _{c}$ and $N$.
Inserting Eq. (\ref{Butterworth2}) into Eqs. (\ref{Analog_Spec_3_low}) and (%
\ref{Analog_Spec_4_low}), we get 
\begin{eqnarray}
-A_{1} &\leq &-10\,\log _{10}\left[ 1+\left( \frac{\Omega }{\Omega _{c}}%
\right) ^{2\,N}\right] \leq 0,\text{ \ \ \ \ }0\leq \Omega \leq \Omega _{1},
\label{Butter_Analog_Pass} \\
-\infty &<&-10\,\log _{10}\left[ 1+\left( \frac{\Omega }{\Omega _{c}}\right)
^{2\,N}\right] \leq -A_{2},\text{ \ \ \ \ }\Omega _{2}\leq \Omega \leq
\infty ,  \label{Butter_Analog_Stop}
\end{eqnarray}
Observe that the right side of Eq. (\ref{Butter_Analog_Pass}) and left side
of Eq. (\ref{Butter_Analog_Stop}) are automatically satisfied. Solving the
left side of Eq. (\ref{Butter_Analog_Pass}) and the right side of Eq. (\ref%
{Butter_Analog_Stop}) for design parameters $N$ and $\Omega _{c}$ we obtain 
\begin{eqnarray}
N &=&\frac{\log _{10}\left( C_{2}/C_{1}\right) }{2\,\log _{10}\left( \Omega
_{2}/\Omega _{1}\right) },  \label{N_Butter} \\
\Omega _{c} &=&\Omega _{1}\,C_{1}^{-1/\left( 2N\right) },
\label{Omega_C_Butter} \\
C_{1} &=&10^{A_{1}/10}-1, \\
C_{2} &=&10^{A_{2}/10}-1.
\end{eqnarray}

\begin{remark}
Filter order $N$ must be chosen as an even integer greater than or equal to
the value given by Eq. (\ref{N_Butter}). With the choice of $\Omega _{c}$
given by Eq. (\ref{Omega_C_Butter}), we satisfy the pass-band requirement of
Eq. (\ref{Analog_Spec_3_low}), while the stop-band requirement of Eq. (\ref%
{Analog_Spec_4_low}) is exceeded. This comes from the choice of $N$ which is
greater than the value specified by Eq. (\ref{N_Butter}).
\end{remark}

\begin{remark}
Because of the analogy mentioned at the end of Section \ref{Introduction},
the parameters of a high-pass Butterworth filter may be derived by first
transforming it to a corresponding low-pass filter and then calculate $N$
and $\Omega _{c}$. The order is the same for the high-pass filter, while $%
\Omega _{c}$ found in this way must be mapped back. Specifically, the
Butterworth parameters of a high-pass filter may be derived in the following
order: Transform the limiting frequencies such that 
\begin{eqnarray}
\Omega _{1,\,lp} &=&\pi -\Omega _{1,\,hp}, \\
\Omega _{2,\,lp} &=&\pi -\Omega _{1,\,hp}, \\
A_{1,\,lp} &=&A_{2,\,hp}, \\
A_{2,\,lp} &=&A_{1,\,hp}.
\end{eqnarray}
Then calculate $N$ and $\Omega _{c,\,lp}$ as explained above. Finally
determine 
\begin{equation}
\Omega _{c,\,hp}=\pi -\Omega _{c,\,lp}.
\end{equation}
\end{remark}

\begin{remark}
Our knowledge of determining the low-pass and high-pass filter parameters
may now be applied to the band-pass and band-stop parameters determination.
The process may be summarized as follows: Divide the specification to
corresponding high-pass and low-pass requirements, as explained at the end
of Section \ref{Introduction}. Then obtain the corresponding filter
parameters as explained above. This gives two even filter orders $N_{1}$ and 
$N_{2}$ as well as two cut-off frequencies $\Omega _{c_{1}}$and $\Omega
_{c_{2}}$. Band-pass or band-stop filter parameters are given by 
\begin{eqnarray}
N &=&2\times \max \left( N_{1},\,N_{2}\right) ,  \label{N_Butter_BP_BS} \\
\Omega _{c} &=&\left[ \Omega _{c_{1}},\,\Omega _{c_{2}}\right] .
\label{Omega_c_Butter_BP_BS}
\end{eqnarray}
Note that we calculate $\Omega _{c_{1}}$ and $\Omega _{c_{2}}$ based on an
even integer found for the low-pass or high-pass requirements. Since the
actual order is twice the maximum of $N_{1}$ and $N_{2}$, both the pass-band
and the stop-band requirements will be exceeded.
\end{remark}

\subsubsection{CHEBYSHEV-1}

An analog, low-pass, Chebyshev-1 filter has parameters $N$ and $\epsilon $.
Inserting Eq. (\ref{Cheb1_Low_Filter}) into Eqs. (\ref{Analog_Spec_3_low})
and (\ref{Analog_Spec_4_low}) we get 
\begin{gather}
-A_{1}\leq -10\,\log _{10}\left[ 1+\epsilon ^{2}T_{N}^{2}\left( \Omega
/\Omega _{1}\right) \right] \leq 0,\text{ \ \ \ \ }0\leq \Omega \leq \Omega
_{1},  \label{Cheb_Analog_Pass} \\
-\infty <-10\,\log _{10}\left[ 1+\epsilon ^{2}T_{N}^{2}\left( \Omega /\Omega
_{1}\right) \right] \leq -A_{2},\text{ \ \ \ \ }\Omega _{2}\leq \Omega
<\infty .  \label{Cheb_Analog_Stop}
\end{gather}
Observe from the properties of Chebyshev polynomials, Eqs. (\ref{ChebProp1})
and (\ref{ChebProp2}), that the right side of Eq. (\ref{Cheb_Analog_Pass})
and the left side of Eq. (\ref{Cheb_Analog_Stop}) are automatically
satisfied. Solving the left side of Eq. (\ref{Cheb_Analog_Pass}) and the
right side of Eq. (\ref{Cheb_Analog_Stop}) for $N$ and $\epsilon $, we
obtain 
\begin{eqnarray}
\epsilon &=&C_{1}^{1/2},  \label{epsilon_Cheb} \\
N &=&\frac{\cosh ^{-1}\left( C_{2}^{1/2}/\epsilon \right) }{\cosh
^{-1}\left( \Omega _{2}/\Omega _{1}\right) },  \label{N_Low_Cheb1}
\end{eqnarray}
where 
\begin{eqnarray}
C_{1} &=&10^{A_{1}/10}-1, \\
\cosh ^{-1}x &=&\ln \left( x+\sqrt{x^{2}-1}\right) .
\end{eqnarray}

\begin{remark}
Again, $N$ must be chosen as the nearest even integer greater than or equal
to the value given by Eq. (\ref{N_Low_Cheb1}). With the choice of $\epsilon $
given by Eq. (\ref{epsilon_Cheb}) the pass-band requirement of Eq. (\ref%
{Analog_Spec_3_low}) is exactly satisfied, while the stop-band requirement
of Eq. (\ref{Analog_Spec_4_low}) is exceeded.
\end{remark}

\begin{remark}
There is a subtle difference between Butterworth and Chebyshev-1 analog
low-pass filters. While cut-off frequency $\Omega _{c},$ in a Butterworth
filter must be calculated from the specifications in Eq. (\ref%
{Omega_C_Butter}), its counterpart $\Omega _{1},$ in a Chebyshev-1 filter is
directly found by frequency prewarping of Eq. (\ref{Dig_2_Analog_Freq}).
\end{remark}

\begin{remark}
Referring to the end of section \ref{Introduction} again, observe that the
mapping of a low-pass to a high-pass Chebyshev-1 does not change $N$ and $%
\epsilon $. Thus the parameters of a digital high-pass filter may be derived
as following: First, set 
\begin{eqnarray}
\Omega _{1,\,lp} &=&\pi -\Omega _{1,\,hp}, \\
\Omega _{2,\,lp} &=&\pi -\Omega _{2,\,hp}, \\
A_{1,\,lp} &=&A_{2,\,hp}, \\
A_{2,\,lp} &=&A_{1,\,hp}.
\end{eqnarray}
Then Determine $N$ and $\epsilon $ as explained above.
\end{remark}

\begin{remark}
Again, we use our knowledge of determining the low-pass and high-pass filter
parameters for the band-pass and band-stop cases. The process is as follows:
Divide the specification to corresponding high-pass and low-pass
requirements, as explained at the end of Section \ref{Introduction}. Then
obtain the corresponding filter parameters for each case as explained above.
This gives two even filter orders $N_{1}$ and $N_{2}$ as well as $%
\varepsilon _{1}$and $\varepsilon _{2}$. Band-pass and band-stop filter
parameters are then obtained by 
\begin{eqnarray}
N &=&2\times \max \left( N_{1},\,N_{2}\right) ,  \label{N_Cheb1_BP_BS} \\
\varepsilon &=&\min \left( \varepsilon _{1},\text{ }\varepsilon _{2}\right) .
\label{Epsilon_Cheb1_BP_BS}
\end{eqnarray}
Since the cut-off frequencies for a band-pass and band-stop Chebyshev filter
are found directly from the filter specifications, Eq. (\ref{N_Cheb1_BP_BS})
has no effect in the pass-band and stop-band values. Thus, as it was the
case for the low-pass and high-pass Chebyshev-1 filters, the pass-band
requirements are exactly satisfied, while the stop-band requirements will be
exceeded.
\end{remark}

\subsection{ANALOG TRANSFER FUNCTION}

To determine the analog, low-pass transfer function, we note from Eqs. (\ref%
{Butter_Real_Part}) and (\ref{Butter_Imag_Part}) that the $2N$ poles of $%
\left| H_{a}\left( s\right) \right| ^{2}$ are complex conjugate pairs and
may be separated to $N$ poles for $H_{a}\left( \sigma +j\,\omega \right) $
and $N$ poles for $H_{a}\left( \sigma -j\,\omega \right) $. Since we are
interested in a stable transfer function, we may choose those poles with
negative real parts, i.e. the poles of \ $H_{a}\left( s\right) $ are found
from 
\begin{equation}
s_{k}=\sigma _{k}+j\,\Omega _{k},\text{ \ \ \ \ }k=0,\,1\,,\,\ldots ,\,N-1,
\end{equation}
where $\sigma _{k}$ and $\Omega _{k}$ are calculated from Eqs. (\ref%
{Butter_Real_Part}), (\ref{Butter_Imag_Part}) for a Butterworth low-pass
filter and from Eqs. (\ref{Cheb_Real_Part}) through (\ref{Gamma}) for a
Chebyshev-1 low-pass filter. Suppose that we have chosen poles $s_{k}$ for $%
k=0,\,1,$\thinspace $\ldots ,\,N/2-1$ such that the complex parts $\Omega
_{k}$ have the same sign. Then $H_{a}\left( s\right) $ may be thought of as
a cascade of $N/2$ transfer functions each of order $2$, so that 
\begin{equation}
H_{a}\left( s\right) =\prod_{k=0}^{N/2-1}\frac{d_{k}}{s^{2}+c_{k}s+d_{k}},
\label{Cascade_Analog}
\end{equation}
where $c_{k}$ and $d_{k}$ are real. This is also a numerically advantageous
approach because it reduces the error associated with the value of
coefficients.

\subsection{DIGITAL SYSTEM FUNCTION}

Each second order factor in Eq. (\ref{Cascade_Analog}) is obtained by
multiplying two complex conjugate poles together. Suppose that the complex
conjugate pair 
\begin{eqnarray}
s_{k} &=&\sigma _{k}+j\,\Omega , \\
s_{k}^{\ast } &=&\sigma _{k}-j\,\Omega ,
\end{eqnarray}
are poles of a transfer function. Then the normalized transfer function may
be expressed as 
\begin{gather}
H_{k}(s)=\frac{r_{k}^{2}}{\left[ s-\left( \sigma _{k}+j\,\Omega _{k}\right) %
\right] \,\left[ s-\left( \sigma _{k}-j\,\Omega _{k}\right) \right] \,}=%
\frac{r_{k}^{2}}{s^{2}-2\,\sigma _{k}s+r_{k}^{2}},  \label{H_2_dim_analog} \\
r_{k}^{2}=\sigma _{k}^{2}+\Omega _{k}^{2}.
\end{gather}
To determine the corresponding system function by the Bilinear
Transformation we substitute $s=\alpha \frac{1-z^{-1}}{1+z^{-1}}$ into Eq. (%
\ref{H_2_dim_analog}) and simplify to get 
\begin{gather}
H\left( z\right) =\frac{x_{0}+x_{1}z^{-1}+x_{2}z^{-2}}{1-\left(
y_{1}z^{-1}+y_{2}z^{-2}\right) },  \label{System_Func_2d} \\
x_{0}=\frac{1}{2}x_{1}=x_{2}=\frac{1}{E}, \\
y_{1}=\frac{2\,r_{k}^{2}-2\,\alpha ^{2}}{E}, \\
y_{2}=-\frac{\alpha ^{2}+2\,\alpha \,\sigma _{k}+r_{k}^{2}}{E}, \\
E=\alpha ^{2}-2\,\sigma _{k}+r_{k}^{2}.
\end{gather}

\subsubsection{FREQUENCY TRANSFORMATIONS}

The digital filter with Eq. (\ref{System_Func_2d}) as its system function is
a low-pass filter with cut-off frequency at $\omega _{c}^{^{\prime }}=1$. To
design a low-pass or high-pass filter with another cut-off frequency, we
need two different kind of mappings, i.e. low-pass to low-pass and low-pass
to high-pass frequency transformations:

\begin{itemize}
\item \textit{Low-pass to low-pass mapping} 
\begin{eqnarray}
Z^{-1} &\rightarrow &\frac{z^{-1}-\beta }{1-\beta \,z^{-1}},
\label{Z_low2low} \\
\beta &=&\frac{\sin \left( \frac{1}{2}-\frac{\omega _{c}}{2}\right) }{\sin
\left( \frac{1}{2}+\frac{\omega _{c}}{2}\right) }.  \label{Beta_low2low}
\end{eqnarray}
The new coefficients may be obtained by 
\begin{gather}
H_{lp}\left( z^{-1}\right) =\frac{\Sigma _{k=0}^{2}a_{k}z^{-k}}{1-\Sigma
_{k=1}^{2}b_{k}z^{-k}}, \\
a_{0}=\frac{x_{0}-x_{1}\beta +x_{2}\beta ^{2}}{D},  \label{a0_Low_Pass} \\
a_{1}=\frac{-2\,x_{0}\beta +x_{1}+x_{1}\beta ^{2}-2\,x_{2}\beta }{D},
\label{a1_Low_Pass} \\
a_{2}=\frac{x_{0}\beta ^{2}-x_{1}\beta +x_{2}}{D}  \label{a2_Low_Pass} \\
b_{1}=\frac{2\,\beta +y_{1}+y_{1}\beta ^{2}-2\,y_{2}\beta }{D},
\label{b1_Low_Pass} \\
b_{2}=\frac{-\beta ^{2}-y_{1}\beta +y_{2}}{D},  \label{b2_Low_Pass} \\
D=1+y_{1}\beta +y_{2}\beta ^{2}.  \label{D_Low_Pass}
\end{gather}

\item \textit{Low-pass to high-pass mapping} 
\begin{eqnarray}
Z^{-1} &\rightarrow &\frac{-z^{-1}-\beta }{1+\beta \,z^{-1}},
\label{Z_low2high} \\
\beta &=&-\frac{\cos \left( \frac{1}{2}+\frac{\omega _{c}}{2}\right) }{\cos
\left( \frac{1}{2}-\frac{\omega _{c}}{2}\right) }.  \label{Beta_low2high}
\end{eqnarray}
The new coefficients may be obtained by 
\begin{gather}
H_{hp}\left( z^{-1}\right) =\frac{\Sigma _{k=0}^{2}a_{k}z^{-k}}{1-\Sigma
_{k=1}^{2}b_{k}z^{-k}},  \label{a0_High_Pass} \\
a_{0}=\frac{x_{0}-x_{1}\beta +x_{2}\beta ^{2}}{D}, \\
a_{1}=-\frac{-2\,x_{0}\beta +x_{1}+x_{1}\beta ^{2}-2\,x_{2}\beta }{D},
\label{a1_High_Pass} \\
a_{2}=\frac{x_{0}\beta ^{2}-x_{1}\beta +x_{2}}{D}  \label{a2_High_Pass} \\
b_{1}=-\frac{2\,\beta +y_{1}+y_{1}\beta ^{2}-2\,y_{2}\beta }{D},
\label{b1_High_Pass} \\
b_{2}=\frac{-\beta ^{2}-y_{1}\beta +y_{2}}{D},  \label{b2_High_Pass} \\
D=1+y_{1}\beta +y_{2}\beta ^{2},  \label{D_High_Pass}
\end{gather}
where $\omega _{c}$ in both cases is the desired digital cut-off frequency.

\item \textit{Low-pass to band-pass mapping} 
\begin{eqnarray}
Z^{-1} &=&-\frac{z^{-2}+\delta \,z^{-1}+\beta }{\beta \,z^{-2}+\delta
\,z^{-1}+1},  \label{Z_low2BandPass} \\
\beta &=&\frac{\zeta -1}{\zeta +1},  \label{Beta_low2BandPass} \\
\delta &=&-\frac{2\,\eta \,\zeta }{\zeta +1},  \label{Eta_low2BandPass} \\
\eta &=&\frac{\cos \left( \frac{\omega _{u}+\omega _{l}}{2}\right) }{\cos
\left( \frac{\omega _{u}-\omega _{l}}{2}\right) },  \label{Ksi_low2BandPass}
\\
\zeta &=&\cot \left( \frac{\omega _{u}-\omega _{l}}{2}\right) \tan \left( 
\frac{1}{2}\right) ,
\end{eqnarray}
where $\omega _{u}$ and $\omega _{l}$ are upper and lower cut-off
frequencies. The new coefficients may be obtained by 
\begin{gather}
H_{bp}\left( z^{-1}\right) =\frac{\Sigma _{k=0}^{4}a_{k}z^{-k}}{1-\Sigma
_{k=1}^{4}b_{k}z^{-k}}, \\
a_{0}=\frac{x_{0}-\beta \,x_{1}+\beta ^{2}x_{2}}{D},  \label{a0_Band_Pass} \\
a_{1}=\frac{2\,\delta \,x_{0}-\left( 1+\beta \right) \,\delta
\,x_{1}+2\,\beta \,\delta \,x_{2}}{D},  \label{a1_Band_Pass} \\
a_{2}=\frac{\left( 2\,\beta +\delta ^{2}\right) \left( x_{0}+x_{2}\right)
-\left( 1+\beta ^{2}+\delta ^{2}\right) x_{1}}{D},  \label{a2_Band_Pass} \\
a_{3}=\frac{2\,\beta \,\delta \,x_{0}-\left( 1+\beta \right) \,\delta
\,x_{1}+2\,\delta \,x_{2}}{D},  \label{a3_Band_Pass} \\
a_{4}=\frac{\beta ^{2}\,x_{0}-\beta \,x_{1}+x_{2}}{D},  \label{a4_Band_Pass}
\\
\,b_{1}=-\frac{2\,\delta +\left( 1+\beta \right) \,\delta \,y_{1}-2\,\beta
\,\delta \,\,y_{2}}{D},  \label{b1_Band_Pass} \\
b_{2}=-\frac{\left( 2\,\beta +\delta ^{2}\right) \left( 1-y_{2}\right)
+\left( 1+\beta ^{2}+\delta ^{2}\right) y_{1}}{D},  \label{b2_Band_Pass} \\
b_{3}=-\frac{2\,\beta \,\delta +\left( 1+\beta \right) \,\delta
\,y_{1}-2\,\delta \,y_{2}}{D},  \label{b3_Band_Pass} \\
b_{4}=-\frac{\beta ^{2}+\beta \,y_{1}-y_{2}}{D},  \label{b4_Band_Pass} \\
D=1+\beta \,y_{1}-\beta ^{2}y_{2}.  \label{D_Band_Pass}
\end{gather}

\item \textit{Low-pass to band-stop mapping} 
\begin{eqnarray}
Z^{-1} &=&\frac{z^{-2}+\delta \,z^{-1}+\beta }{\beta \,z^{-2}+\delta
\,z^{-1}+1},  \label{Z_low2BandStop} \\
\beta &=&\frac{1-\zeta }{1+\zeta },  \label{Beta_low2BandStop} \\
\delta &=&-\frac{2\,\eta }{1+\zeta },  \label{Delta_low2BandStop} \\
\eta &=&\frac{\cos \left( \frac{\omega _{u}+\omega _{l}}{2}\right) }{\cos
\left( \frac{\omega _{u}-\omega _{l}}{2}\right) },  \label{Eta_low2BandStop}
\\
\zeta &=&\tan \left( \frac{\omega _{u}-\omega _{l}}{2}\right) \tan \left( 
\frac{1}{2}\right) .  \label{Ksi_low2BandStop}
\end{eqnarray}
where $\omega _{u}$ and $\omega _{l}$ are upper and lower cut-off
frequencies. The new coefficients may be obtained by 
\begin{gather}
H_{bp}\left( z^{-1}\right) =\frac{\Sigma _{k=0}^{4}a_{k}z^{-k}}{1-\Sigma
_{k=1}^{4}b_{k}z^{-k}}, \\
a_{0}=\frac{x_{0}+\beta \,x_{1}+\beta ^{2}x_{2}}{D},  \label{a0_Band_Stop} \\
a_{1}=\frac{2\,\delta \,x_{0}+\left( 1+\beta \right) \,\delta
\,x_{1}+2\,\beta \,\delta \,x_{2}}{D},  \label{a1_Band_Stop} \\
a_{2}=\frac{\left( 2\,\beta +\delta ^{2}\right) \left( x_{0}+x_{2}\right)
+\left( 1+\beta ^{2}+\delta ^{2}\right) x_{1}}{D},  \label{a2_Band_Stop} \\
a_{3}=\frac{2\,\beta \,\delta \,x_{0}+\left( 1+\beta \right) \,\delta
\,x_{1}+2\,\delta \,x_{2}}{D},  \label{a3_Band_Stop} \\
a_{4}=\frac{\beta ^{2}\,x_{0}+\beta \,x_{1}+x_{2}}{D},  \label{a4_Band_Stop}
\\
\,b_{1}=-\frac{2\,\delta -\left( 1+\beta \right) \,\delta \,y_{1}-2\,\beta
\,\delta \,\,y_{2}}{D},  \label{b1_Band_Stop} \\
b_{2}=-\frac{\left( 2\,\beta +\delta ^{2}\right) \left( 1-y_{2}\right)
-\left( 1+\beta ^{2}+\delta ^{2}\right) y_{1}}{D},  \label{b2_Band_Stop} \\
b_{3}=-\frac{2\,\beta \,\delta -\left( 1+\beta \right) \,\delta
\,y_{1}-2\,\delta \,y_{2}}{D},  \label{b3_Band_Stop} \\
b_{4}=-\frac{\beta ^{2}-\beta \,y_{1}-y_{2}}{D},  \label{b4_Band_Stop} \\
D=1-\beta \,y_{1}-\beta ^{2}y_{2}.  \label{D_Band_Stop}
\end{gather}
\end{itemize}

\begin{remark}
The transformation given by Eqs. (\ref{Z_low2low}) just changes the
frequency from $1$ to $\omega _{c}$, while the frequency mapping of Eq. (\ref%
{Beta_low2high}), (\ref{Beta_low2BandPass}) or (\ref{Beta_low2BandStop}) not
only changes the filter property from low-pass to high-pass, band-pass or
band-stop, but also the cut-off frequency from $1$ to $\omega _{c}$ for a
high-pass and from 1 to $\omega _{u}$ and $\omega _{l}$ for a band-pass or
band-stop filter.
\end{remark}

\begin{remark}
By comparing Eqs. (\ref{a0_Low_Pass}) through (\ref{D_Low_Pass}) with Eqs. (%
\ref{a0_High_Pass}) through (\ref{D_High_Pass}), we see that by applying the
correct value of $\beta $, high-pass coefficients may be obtained just by
changing the signs of low-pass coefficients $a_{1}$ and $b_{1}$.
\end{remark}

\begin{remark}
By comparing Eqs. (\ref{Z_low2low}) and (\ref{Z_low2high}) with (\ref%
{Z_low2BandPass}) and (\ref{Z_low2BandStop}), it is clear that the first two
transformations keep the filter order, while the last two doubles it. This
is the reason for doubling maximum $N$ to determine the band-pass and
band-stop filter order in Eq. (\ref{N_Butter_BP_BS}) and (\ref{N_Cheb1_BP_BS}%
).
\end{remark}

Frequency transformations are also discussed in \cite{Smith} pages 628-630
and in \cite{Mitra}, section 5-6.

\subsubsection{CASCADING SYSTEM FUNCTIONS}

Remember that we calculate coefficients of a second order system function in
each step and then put it in series with the result of previous system by
integrating the new coefficients into global vectors $\vec{a}$ and $\vec{b}$%
. Suppose now that two second order system functions are defined by 
\begin{eqnarray}
H_{1}\left( z\right) &=&\frac{\sum_{k=0}^{4}\alpha _{k}z^{-k}}{%
1-\sum_{k=1}^{4}\beta _{k}z^{-k}},  \label{H_1} \\
H_{2}\left( z\right) &=&\frac{\sum_{k=0}^{n}\delta _{k}z^{-k}}{%
1-\sum_{k=1}^{n}\eta _{k}z^{-k}},  \label{H_2}
\end{eqnarray}
and define the third system function by cascading $H_{1}\left( z\right) $
and $H_{2}\left( z\right) $ such that 
\begin{equation}
H\left( z\right) =H_{1}\left( z\right) \times H_{2}\left( z\right) =\frac{%
\sum_{k=0}^{n+4}c_{k}z^{-k}}{1-\sum_{k=1}^{n+2}d_{k}z^{-k}}.
\label{Cascade_Digital}
\end{equation}

Fortunately, multiplication of two polynomials is accomplished efficiently
just by convolving their corresponding coefficients. Define now two
polynomials of orders $m$ and $n$ by $P_{1}=\sum_{k=0}^{m}\theta _{k}z^{-k}$
and $P_{2}=\sum_{k=0}^{n}\lambda _{k}z^{-k}$. Then, it may be verified that
the coefficients of polynomial $P_{3}=P_{1}\times P_{2}=$ $%
\sum_{k=0}^{m+n}\varphi _{k}z^{-k}$ are obtained by 
\begin{eqnarray}
\vec{\varphi} &=&\vec{\theta}\ast \vec{\lambda}, \\
\varphi _{k} &=&\sum_{i=0}^{m+n}\theta _{i}\lambda _{k-i},\text{ \ \ \ \ }%
k=0,\,1,\,\ldots ,\,m+n.
\end{eqnarray}

This result can be applied to the coefficients of the two system functions.
Define now $\vec{\alpha}=\left[ a_{0},\,a_{1},\,a_{2},\,a_{3},\,a_{4}\right] 
$, $\vec{\beta}=\left[ 1,\,-b_{1},\,-b_{2},\,-b_{3},\,-b_{4}\right] $, $\vec{%
\delta}=\left[ A_{0},\,\,A_{1},\,\ldots ,\,\,A_{n}\right] $, $\vec{\eta}=%
\left[ 1,\,-B_{1},\,\ldots ,\,-B_{n}\right] $. Then the coefficients of $%
H\left( z\right) $ in Eq. (\ref{Cascade_Digital}) are obtained by 
\begin{eqnarray}
c_{k} &=&\sum_{i=0}^{n+4}\alpha _{i}\delta _{k-i},\text{ \ \ \ \ }%
k=0,\,1,\,\,\ldots ,\,n+4, \\
d_{k} &=&\sum_{i=0}^{n+4}\beta _{i}\eta _{k-i},\text{ \ \ \ \ }%
k=1,\,\,\ldots ,\,n+2.
\end{eqnarray}

\begin{remark}
For a low-pass or a high-pass filter the numerator and denominator of $%
H_{1}\left( z\right) $ in Eq. (\ref{H_1}) are of order 2, thus $%
a_{3}=a_{4}=b_{3}=b_{4}=0.$
\end{remark}

\begin{remark}
Since we chose a stable analog prototype and the Bilinear Transformation
preserves the stability property, we are guaranteed the
unconditional-stability of the resultant digital filter.
\end{remark}

\subsection{AMPLITUDE NORMALIZATION}

By now we have found the frequency-response of the form 
\begin{equation}
H\left( z=e^{j\,\omega }\right) =\frac{\sum_{k=0}^{N}a_{k}e^{-j\,k\,\omega }%
}{1-\sum_{k=1}^{N}b_{k}e^{-j\,k\,\omega }}.  \label{H_Non_Normal}
\end{equation}
Initially, the absolute value of $H$ at $\omega =\omega _{c}$ may be
obtained by 
\begin{equation}
A_{initial}=\left| H\left( \omega =\omega _{c}\right) \right| =\left| \frac{%
\sum_{k=0}^{N}a_{k}e^{-j\,k\,\omega _{c}}}{1-\sum_{k=1}^{N}b_{k}e^{-j\,k\,%
\omega _{c}}}\right| .  \label{A_initial}
\end{equation}
It is sensible to normalize Eq. (\ref{H_Non_Normal}) such that $H\left(
\omega =\omega _{c}\right) =$ $H_{a}\left( \Omega =\Omega _{c}\right) $.
Denote by $A_{desired}$ the desired amplitude for a digital filter at $%
\omega =\omega _{c}$ and note that for an even $N$%
\begin{equation}
A_{desired}=\left| H_{a}\left( \Omega =\Omega _{c}\right) \right| =\left\{ 
\begin{array}{cc}
1, & \text{for a Butterworth filter} \\ 
\frac{1}{\sqrt{1+\epsilon ^{2}}}, & \text{for a Chebyshev-1 filter}%
\end{array}
\right.  \label{A_desired}
\end{equation}
The amplitude may now be normalized by 
\begin{equation}
H_{nornal}=H\left( e^{j\,\omega }\right) \times \frac{A_{desired}}{%
A_{initial}}.  \label{Amp_normal}
\end{equation}
Inserting Eq. (\ref{H_Non_Normal}) into (\ref{Amp_normal}), this is
equivalent to normalizing the numerator coefficients in Eq. (\ref%
{H_Non_Normal}). Thus, the normalized coefficients of the digital filter are
calculated from 
\begin{eqnarray}
\vec{a}_{normal} &=&\vec{a}\times \frac{A_{desired}}{A_{initial}}, \\
\vec{b}_{normal} &=&\vec{b},
\end{eqnarray}
where $A_{initial}$ and $A_{desired}$ are obtained from Eqs. (\ref{A_initial}%
) and (\ref{A_desired}), while $\vec{a}$ and $\vec{b}$ are the previous
coefficient vectors.

\subsection{IMPULSE-RESPONSE CONSTRUCTION}

Having found the digital filter coefficients, we may obtain the
impulse-response in one of the two ways:

\begin{itemize}
\item \textit{Direct or exact method}: Calculate $H\left( e^{j\,\omega
}\right) $ by inserting $\vec{a}$ and $\vec{b}$ into Eq. (\ref{Freq_Resp}),
then apply $\mathcal{F}^{-1}$-transform to $H\left( e^{j\,\omega }\right) $, 
\begin{equation}
h\left[ n\right] =\mathcal{F}^{-1}\left\{ H\left( e^{j\,\omega }\right)
\right\} .
\end{equation}

\item \textit{Backward method}: Send an impulse into the system defined by
difference Eq. (\ref{LTI_Causal_Filter}) and calculate the output, 
\begin{equation}
x\left[ n\right] =\delta \left[ n\right] \Leftrightarrow y\left[ n\right] =h%
\left[ n\right] .
\end{equation}
\end{itemize}

Here, we apply the second method to construct the $h\left[ n\right] $. Then
we compare its magnitude response with the results of Direct Calculations.

\begin{remark}
Clearly, to calculate $h\left[ n\right] $ in the second method, we have to
send a finite length impulse. As we see later this truncation is an error
source which can be seen when we compare the magnitude responses of the two
methods. Intuitively, we expect smaller deviation when impulse-response
length increases.
\end{remark}

\subsection{SCALING FILTER COEFFICIENTS}

In some cases we are just able to handle integer coefficients, so we must
find a way to convert all coefficients to integers. Suppose now that $L$ is
the largest representable integer in the actual processor and $C$ is some
unknown constant. Denote by $a_{k}$ and $b_{k}$ the normalized coefficients
obtained above and define a new filter by 
\begin{eqnarray}
\acute{a}_{k} &=&C\,a_{k},\;\;\;\;\;\text{for }k=0,\,1,\,2,\ldots ,\,N,
\label{Scaled_a} \\
\acute{b}_{k} &=&C\,b_{k},\;\;\;\;\;\text{for }k=0,\,1,\,2,\ldots ,\,N,
\label{Scaled_B} \\
\tilde{y}\left[ n\right] &=&\sum_{k=0}^{N}\acute{a}_{k}x\left[ n-k\right]
+\sum_{k=1}^{N}\acute{b}_{k}\tilde{y}\left[ n-k\right] ,\text{ \ \ \ \ }%
\forall \,n,  \label{Scaled_y} \\
\acute{y}\left[ n\right] &=&\frac{1}{C}\tilde{y}\left[ n\right] ,\text{\ \ \
\ \ }\forall \,n\text{.}  \label{Correct_y}
\end{eqnarray}
It may be verified that $\acute{y}\left[ n\right] $ found in Eq. (\ref%
{Correct_y}) is the same as the output of Eq. (\ref{LTI_Causal_Filter}). Our
goal is now to determine scaling factor $C$. As we see from Eq. (\ref%
{Correct_y}), we must divide the output found from Eq. (\ref{Scaled_y}) by $%
C $ in each step. Since the original decimal coefficients $a_{k}$ and $b_{k}$
may be quite small, we first divide them by the minimum positive magnitude
coefficient larger than a predefined value $\delta $. Then multiply all new
coefficients by a positive integer such that the largest coefficient
obtained is still inside the integer range. Thus 
\begin{eqnarray}
m_{1} &=&\max \left[ \min_{k=0,\,1,\,\ldots ,\,N}abs\left( a_{k}\right)
,\,\delta \right] ,  \label{Min_a_k} \\
m_{2} &=&\max \left[ \min_{k=1,\,\ldots ,\,N}abs\left( b_{k}\right)
,\,\delta \right] ,  \label{Min_b_k} \\
m_{3} &=&\min \left( m_{1},\,m_{2}\right) ,  \label{Min_ab} \\
M_{1} &=&\max_{k=0,\,1,\,\ldots ,\,N}\left[ abs\left( \frac{a_{k}}{m_{3}}%
\right) \right] ,  \label{Max_a_k_plus} \\
M_{2} &=&\max_{k=1,\,\ldots ,\,N}\left[ abs\left( \frac{b_{k}}{m_{3}}\right) %
\right] ,  \label{Max_b_k_plus} \\
M_{3} &=&\max \left( M_{1},\,M_{2}\right) ,  \label{Max_ab_plus} \\
C &=&floor\left( \frac{L}{m_{3}\times M_{3}}\right) ,  \label{Scale_Factor}
\\
\acute{a}_{k} &=&floor\left( C\times a_{k}\right)
,\;\;\;\;\;k=0,\,1,\,\ldots ,\,N,  \label{A_r} \\
\acute{b}_{k} &=&floor\left( C\times b_{k}\right) ,\;\;\;\;\;k=1,\,\ldots
,\,N.  \label{B_r}
\end{eqnarray}

\begin{remark}
Note that function $floor$ return an integer by truncating all digits after
the decimal point. The new coefficients obtained from Eqs. (\ref{A_r}) and (%
\ref{B_r}) are used in the Backward Method to calculate the
impulse-response. Thus in the Backward Calculations we have two sources of
numeric error, i.e. the truncation of the impulse-response and the rounding
effect introduced by Eq. (\ref{Scale_Factor}).
\end{remark}

\begin{remark}
We have done more than just finding the magnitude of the smallest
coefficient in Eqs. (\ref{Min_a_k}) and (\ref{Min_b_k}). The reason is the
case where dividing by the smallest coefficient is enough to go outside the
integer range of the actual processor. Furthermore we may do all these
calculations in a while-loop, where we increase $\delta $ to obtain a value
for the minimum coefficient, which after division gives a maximum
coefficient inside the range. If it still is impossible keep the integers
inside the range we have no choice than stop calculations and give a warning
message about the problem. In this case we are left with the decimal
coefficients which will be used as described above.
\end{remark}

\section{RESULTS}

In this section we present results of the digital filter implementation.
Because the input data are obtained by sampling a continuous function, it is
more natural to use the normalized frequency with respect to the sampling
frequency, i.e. $0\leq f=\frac{\omega }{2\pi }\leq 0.5.$ To be able to
compare Butterworth with Chebyshev-1 we applied the same input data for
low-pass, high-pass, band-pass and bans-stop filters. Furthermore we will
compare the results of the Direct with the Backward calculations. In the
first method, we find filter coefficients $a_{k}$ and $b_{k}$, then obtain
the frequency-response from $H\left( e^{j\,\omega }\right) =\frac{%
\sum_{k=0}^{N}a_{k}\,e^{-j\,\omega \,k}}{1-\sum_{k=1}^{N}b_{k}\,e^{-j\,%
\omega \,k}}$. In the second method, we begin by reading the data files
containing integer vectors $\acute{a}_{k}$, $\acute{b}_{k}$ and integer
scale factor $C$, apply Eq. (\ref{LTI_Causal_Filter}) with $\delta \left[ n%
\right] $ as the input to construct $h\left[ n\right] $ and finally apply $%
\mathcal{F}$-operator to obtain its frequency-response. Then we plot both
magnitude responses and compare them with each other.

\begin{center}
\FRAME{ftbphFU}{11.2511cm}{8.468cm}{0pt}{\Qcb{Amplitude of a digital
low-pass Butterworth filter with limiting frequencies and attenuations equal
to $f_{\lim }=[0.20$ $\ 0.25]$ $Hz$ and $A_{\lim }=[1\,\ 20]$ $dB$. Backward
calculations are performed on a 16-bits processor.}}{\Qlb{Low_Butt_Res}}{%
low_butt_res.eps}{\special{language "Scientific Word";type
"GRAPHIC";maintain-aspect-ratio TRUE;display "ICON";valid_file "F";width
11.2511cm;height 8.468cm;depth 0pt;original-width 8.003in;original-height
6.0105in;cropleft "0";croptop "1";cropright "1";cropbottom "0";filename
'Low_Butt_Res.eps';file-properties "XNPEU";}}\FRAME{ftbphFU}{11.2467cm}{%
8.468cm}{0pt}{\Qcb{$dB-$amplitude of a digital low-pass Butterworth filter
with limiting frequencies and attenuations equal to $f_{\lim }=[0.20$ $\
0.25]$ $Hz$ and $A_{\lim }=[1\,\ 20]$ $dB$. Backward calculations are
performed on a 16-bits processor.}}{\Qlb{Low_Butt_dB_Res}}{%
low_butt_db_res.eps}{\special{language "Scientific Word";type
"GRAPHIC";maintain-aspect-ratio TRUE;display "ICON";valid_file "F";width
11.2467cm;height 8.468cm;depth 0pt;original-width 8.003in;original-height
6.0105in;cropleft "0";croptop "1.0001";cropright "0.9996";cropbottom
"0";filename 'Low_Butt_dB_Res.eps';file-properties "XNPEU";}}\FRAME{ftbphFU}{%
11.2511cm}{8.468cm}{0pt}{\Qcb{Deviation between direct and backward
calculations of the low-pass Butterworth magnitude on a 16-bits processor.}}{%
\Qlb{Low_Butt_Error}}{low_butt_error.eps}{\special{language "Scientific
Word";type "GRAPHIC";maintain-aspect-ratio TRUE;display "ICON";valid_file
"F";width 11.2511cm;height 8.468cm;depth 0pt;original-width
8.003in;original-height 6.0105in;cropleft "0";croptop "1";cropright
"1";cropbottom "0";filename 'Low_Butt_Error.eps';file-properties "XNPEU";}}%
\FRAME{ftbphFU}{11.2511cm}{8.468cm}{0pt}{\Qcb{Amplitude of a digital
low-pass Chebyshev-1 filter with limiting frequencies and attenuations equal
to $f_{\lim }=[0.20$ $\ 0.25]$ $Hz$ and $A_{\lim }=[1\,\ 20]$ $dB$. Backward
calculations are performed on a 16-bits processor.}}{\Qlb{Low_Cheb1_Res}}{%
low_cheb_res.eps}{\special{language "Scientific Word";type
"GRAPHIC";maintain-aspect-ratio TRUE;display "ICON";valid_file "F";width
11.2511cm;height 8.468cm;depth 0pt;original-width 8.003in;original-height
6.0105in;cropleft "0";croptop "1";cropright "1";cropbottom "0";filename
'Low_Cheb_Res.eps';file-properties "XNPEU";}}\FRAME{ftbphFU}{11.2511cm}{%
8.468cm}{0pt}{\Qcb{$dB-$amplitude of a digital low-pass Chebyshev-1 filter
with limiting frequencies and attenuations equal to $f_{\lim }=[0.20$ $\
0.25]$ $Hz$ and $A_{\lim }=[1\,\ 20]$ $dB$. Backward calculations are
performed on a 16-bits processor.}}{\Qlb{Low_Cheb1_dB_Res}}{%
low_cheb_db_res.eps}{\special{language "Scientific Word";type
"GRAPHIC";maintain-aspect-ratio TRUE;display "ICON";valid_file "F";width
11.2511cm;height 8.468cm;depth 0pt;original-width 8.003in;original-height
6.0105in;cropleft "0";croptop "1";cropright "1";cropbottom "0";filename
'Low_Cheb_dB_Res.eps';file-properties "XNPEU";}}\FRAME{ftbphFU}{11.2511cm}{%
8.468cm}{0pt}{\Qcb{Deviation between direct and backward calculations of the
low-pass Chebyshev-1 magnitude on a 16-bits processor.}}{\Qlb{Low_Cheb1_Error%
}}{low_cheb_error.eps}{\special{language "Scientific Word";type
"GRAPHIC";maintain-aspect-ratio TRUE;display "ICON";valid_file "F";width
11.2511cm;height 8.468cm;depth 0pt;original-width 8.003in;original-height
6.0105in;cropleft "0";croptop "1";cropright "1";cropbottom "0";filename
'Low_Cheb_Error.eps';file-properties "XNPEU";}}\FRAME{ftbphFU}{11.2511cm}{%
8.468cm}{0pt}{\Qcb{Amplitude of a digital high-pass Butterworth filter with
limiting frequencies and attenuations equal to $f_{\lim }=[0.30$ $\ 0.35]$ $%
Hz$ and $A_{\lim }=[40\,\ 1]$ $dB$. Backward calculations are performed on a
32-bits processor.}}{\Qlb{High_Butt_Res}}{high_butt_res.eps}{\special%
{language "Scientific Word";type "GRAPHIC";maintain-aspect-ratio
TRUE;display "ICON";valid_file "F";width 11.2511cm;height 8.468cm;depth
0pt;original-width 8.003in;original-height 6.0105in;cropleft "0";croptop
"1";cropright "1";cropbottom "0";filename
'High_Butt_Res.eps';file-properties "XNPEU";}}\FRAME{ftbphFU}{11.2511cm}{%
8.468cm}{0pt}{\Qcb{$dB-$amplitude of a digital high-pass Butterworth filter
with limiting frequencies and attenuations equal to $f_{\lim }=[0.30$ $\
0.35]$ $Hz$ and $A_{\lim }=[40\,\ 1]$ $dB$. Backward calculations are
performed on a 32-bits processor.}}{\Qlb{High_Butt_dB_Res}}{%
high_butt_db_res.eps}{\special{language "Scientific Word";type
"GRAPHIC";maintain-aspect-ratio TRUE;display "ICON";valid_file "F";width
11.2511cm;height 8.468cm;depth 0pt;original-width 8.003in;original-height
6.0105in;cropleft "0";croptop "1";cropright "1";cropbottom "0";filename
'High_Butt_dB_Res.eps';file-properties "XNPEU";}}\FRAME{ftbphFU}{11.2511cm}{%
8.468cm}{0pt}{\Qcb{Deviation between direct and backward calculations of the
high-pass Butterworth magnitude on a 32-bits processor.}}{\Qlb{%
High_Butt_Error}}{high_butt_error.eps}{\special{language "Scientific
Word";type "GRAPHIC";maintain-aspect-ratio TRUE;display "ICON";valid_file
"F";width 11.2511cm;height 8.468cm;depth 0pt;original-width
8.003in;original-height 6.0105in;cropleft "0";croptop "1";cropright
"1";cropbottom "0";filename 'High_Butt_Error.eps';file-properties "XNPEU";}}%
\FRAME{ftbphFU}{11.2467cm}{8.468cm}{0pt}{\Qcb{Amplitude of a digital
high-pass Chebyshev-1 filter with limiting frequencies and attenuations
equal to $f_{\lim }=[0.30$ $\ 0.35]$ $Hz$ and $A_{\lim }=[40\,\ 1]$ $dB$.
Backward calculations are performed on a 32-bits processor.}}{\Qlb{%
High_Cheb1_Res}}{high_cheb_res.eps}{\special{language "Scientific Word";type
"GRAPHIC";maintain-aspect-ratio TRUE;display "ICON";valid_file "F";width
11.2467cm;height 8.468cm;depth 0pt;original-width 8.003in;original-height
6.0105in;cropleft "0";croptop "1.0001";cropright "0.9996";cropbottom
"0";filename 'High_Cheb_Res.eps';file-properties "XNPEU";}}\FRAME{ftbphFU}{%
11.2511cm}{8.468cm}{0pt}{\Qcb{$dB-$amplitude of a digital high-pass
Chebyshev-1 filter with limiting frequencies and attenuations equal to $%
f_{\lim }=[0.30$ $\ 0.35]$ $Hz$ and $A_{\lim }=[40\,\ 1]$ $dB$. Backward
calculations are performed on a 32-bits processor.}}{\Qlb{High_Cheb1_dB_Res}%
}{high_cheb_db_res.eps}{\special{language "Scientific Word";type
"GRAPHIC";maintain-aspect-ratio TRUE;display "ICON";valid_file "F";width
11.2511cm;height 8.468cm;depth 0pt;original-width 8.003in;original-height
6.0105in;cropleft "0";croptop "1";cropright "1";cropbottom "0";filename
'High_Cheb_dB_Res.eps';file-properties "XNPEU";}}

\FRAME{ftbphFU}{11.2511cm}{8.468cm}{0pt}{\Qcb{Deviation between direct and
backward calculations of the high-pass Chebyshev-1 magnitude on a 32-bits
processor.}}{\Qlb{High_Cheb1_Error}}{high_cheb_error.eps}{\special{language
"Scientific Word";type "GRAPHIC";maintain-aspect-ratio TRUE;display
"ICON";valid_file "F";width 11.2511cm;height 8.468cm;depth
0pt;original-width 8.003in;original-height 6.0105in;cropleft "0";croptop
"1";cropright "1";cropbottom "0";filename
'High_Cheb_Error.eps';file-properties "XNPEU";}}

\FRAME{ftbphFU}{11.2467cm}{8.468cm}{0pt}{\Qcb{Amplitude of a digital
band-pass, Butterworth filter with limiting frequencies and attenuations
equal to $f_{\lim }=[0.20$ $\ 0.25$ $0.35$ $\ 0.4]$ and $A_{\lim }=[20\,\ 1$
\ $1$ \ $30]$. Backward calculations are performed on a 32-bits processor.}}{%
\Qlb{BandPass_Butt_Res}}{bandpass_butt_res.eps}{\special{language
"Scientific Word";type "GRAPHIC";maintain-aspect-ratio TRUE;display
"ICON";valid_file "F";width 11.2467cm;height 8.468cm;depth
0pt;original-width 8.003in;original-height 6.0105in;cropleft "0";croptop
"1";cropright "0.9995";cropbottom "0";filename
'BandPass_Butt_Res.eps';file-properties "XNPEU";}}

%TCIMACRO{\TeXButton{clearpage}{\clearpage}}%
%BeginExpansion
\clearpage%
%EndExpansion

\FRAME{ftbphFU}{11.2467cm}{8.468cm}{0pt}{\Qcb{$dB-$amplitude of a digital
band-pass, Butterworth filter with limiting frequencies and attenuations
equal to $f_{\lim }=[0.20$ $\ 0.25$ $0.35$ $\ 0.4]$ and $A_{\lim }=[20\,\ 1$
\ $1$ \ $30]$. Backward calculations are performed on a 32-bits processor.}}{%
\Qlb{Bandpass_Butt_dB_Res}}{bandpass_butt_db_res.eps}{\special{language
"Scientific Word";type "GRAPHIC";maintain-aspect-ratio TRUE;display
"ICON";valid_file "F";width 11.2467cm;height 8.468cm;depth
0pt;original-width 8.003in;original-height 6.0105in;cropleft "0";croptop
"1.0001";cropright "0.9996";cropbottom "0";filename
'BandPass_Butt_dB_Res.eps';file-properties "XNPEU";}}\FRAME{ftbphFU}{%
11.2467cm}{8.468cm}{0pt}{\Qcb{Deviation between direct and backward
calculations of the band-pass Butterworth magnitude on a 32-bits processor.}%
}{\Qlb{BandPass_Butt_Error}}{bandpass_butt_error.eps}{\special{language
"Scientific Word";type "GRAPHIC";maintain-aspect-ratio TRUE;display
"ICON";valid_file "F";width 11.2467cm;height 8.468cm;depth
0pt;original-width 8.003in;original-height 6.0105in;cropleft "0";croptop
"1.0001";cropright "0.9996";cropbottom "0";filename
'BandPass_Butt_Error.eps';file-properties "XNPEU";}}\FRAME{ftbphFU}{11.2467cm%
}{8.468cm}{0pt}{\Qcb{Amplitude of a digital band-pass, Chebyshev-1 filter
with limiting frequencies and attenuations equal to $f_{\lim }=[0.20$ $\
0.25 $ $0.35$ $\ 0.4]$ and $A_{\lim }=[20\,\ 1$ \ $1$ \ $30]$. Backward
calculations are performed on a 32-bits processor.}}{\Qlb{BandPass_Cheb1_Res}%
}{bandpass_cheb_res.eps}{\special{language "Scientific Word";type
"GRAPHIC";maintain-aspect-ratio TRUE;display "ICON";valid_file "F";width
11.2467cm;height 8.468cm;depth 0pt;original-width 8.003in;original-height
6.0105in;cropleft "0";croptop "1.0001";cropright "0.9996";cropbottom
"0";filename 'BandPass_Cheb_Res.eps';file-properties "XNPEU";}}\FRAME{ftbphFU%
}{11.2511cm}{8.468cm}{0pt}{\Qcb{$dB-$amplitude of a digital band-pass,
Chebyshev-1 filter with limiting frequencies and attenuations equal to $%
f_{\lim }=[0.20$ $\ 0.25$ $0.35$ $\ 0.4]$ and $A_{\lim }=[20\,\ 1$ \ $1$ \ $%
30]$. Backward calculations are performed on a 32-bits processor.}}{\Qlb{%
BandPass_Cheb1_dB_Res}}{bandpass_cheb_db_res.eps}{\special{language
"Scientific Word";type "GRAPHIC";maintain-aspect-ratio TRUE;display
"ICON";valid_file "F";width 11.2511cm;height 8.468cm;depth
0pt;original-width 8.003in;original-height 6.0105in;cropleft "0";croptop
"1";cropright "1";cropbottom "0";filename
'BandPass_Cheb_dB_Res.eps';file-properties "XNPEU";}}\FRAME{ftbphFU}{%
11.2467cm}{8.468cm}{0pt}{\Qcb{Deviation between direct and backward
calculations of the band-pass Chebyshev-1 magnitude on a 32-bits processor.}%
}{\Qlb{BandPass_Cheb1_Error}}{bandpass_cheb_error.eps}{\special{language
"Scientific Word";type "GRAPHIC";maintain-aspect-ratio TRUE;display
"ICON";valid_file "F";width 11.2467cm;height 8.468cm;depth
0pt;original-width 8.003in;original-height 6.0105in;cropleft "0";croptop
"1.0001";cropright "0.9996";cropbottom "0";filename
'BandPass_Cheb_Error.eps';file-properties "XNPEU";}}\FRAME{ftbphFU}{11.2467cm%
}{8.468cm}{0pt}{\Qcb{Amplitude of a digital band-stop, Butterworth filter
with limiting frequencies and attenuations equal to $f_{\lim }=[0.20$ $\ 0.3$
$\ 0.4$ $\ 0.45]$ and $A_{\lim }=[1$ $\ 20\,\ 35$ \ $1]$. Backward
calculations are performed on a 32-bits processor.}}{\Qlb{BandStop_Butt_Res}%
}{bandstop_butt_res.eps}{\special{language "Scientific Word";type
"GRAPHIC";maintain-aspect-ratio TRUE;display "ICON";valid_file "F";width
11.2467cm;height 8.468cm;depth 0pt;original-width 8.003in;original-height
6.0105in;cropleft "0";croptop "1.0001";cropright "0.9996";cropbottom
"0";filename 'BandStop_Butt_Res.eps';file-properties "XNPEU";}}\FRAME{ftbphFU%
}{11.2467cm}{8.468cm}{0pt}{\Qcb{$dB-$amplitude of a digital band-stop,
Butterworth filter with limiting frequencies and attenuations equal to $%
f_{\lim }=[0.20$ $\ 0.3$ $\ 0.4$ $\ 0.45]$ and $A_{\lim }=[1$ $\ 20\,\ 35$ \ 
$1]$. Backward calculations are performed on a 32-bits processor.}}{\Qlb{%
BandStop_Butt_dB_Res}}{bandstop_butt_db_res.eps}{\special{language
"Scientific Word";type "GRAPHIC";maintain-aspect-ratio TRUE;display
"ICON";valid_file "F";width 11.2467cm;height 8.468cm;depth
0pt;original-width 8.003in;original-height 6.0105in;cropleft "0";croptop
"1.0001";cropright "0.9996";cropbottom "0";filename
'BandStop_Butt_dB_Res.eps';file-properties "XNPEU";}}\FRAME{ftbphFU}{%
11.2467cm}{8.468cm}{0pt}{\Qcb{Deviation between direct and backward
calculations of the band-stop Butterworth magnitude on a 32-bits processor.}%
}{\Qlb{BandStop_Butt_Error}}{bandstop_butt_error.eps}{\special{language
"Scientific Word";type "GRAPHIC";maintain-aspect-ratio TRUE;display
"ICON";valid_file "F";width 11.2467cm;height 8.468cm;depth
0pt;original-width 8.003in;original-height 6.0105in;cropleft "0";croptop
"1.0001";cropright "0.9996";cropbottom "0";filename
'BandStop_Butt_Error.eps';file-properties "XNPEU";}}\FRAME{ftbphFU}{11.2511cm%
}{8.468cm}{0pt}{\Qcb{Amplitude of a digital band-stop, Chebyshev-1 filter
with limiting frequencies and attenuations equal to $f_{\lim }=[0.20$ $\ 0.3$
$\ 0.4$ $\ 0.45]$ and $A_{\lim }=[1$ $\ 20\,\ 35$ \ $1]$. Backward
calculations are performed on a 32-bits processor.}}{\Qlb{BandStop_Cheb1_Res}%
}{bandstop_cheb_res.eps}{\special{language "Scientific Word";type
"GRAPHIC";maintain-aspect-ratio TRUE;display "ICON";valid_file "F";width
11.2511cm;height 8.468cm;depth 0pt;original-width 8.003in;original-height
6.0105in;cropleft "0";croptop "1";cropright "1";cropbottom "0";filename
'BandStop_Cheb_Res.eps';file-properties "XNPEU";}}\FRAME{ftbphFU}{11.2467cm}{%
8.468cm}{0pt}{\Qcb{$dB-$amplitude of a digital band-stop, Chebyshev-1 filter
with limiting frequencies and attenuations equal to $f_{\lim }=[0.20$ $\ 0.3$
$\ 0.4$ $\ 0.45]$ and $A_{\lim }=[1$ $\ 20\,\ 35$ \ $1]$. Backward
calculations are performed on a 32-bits processor.}}{\Qlb{%
BandStop_Cheb1_dB_Res}}{bandstop_cheb_db_res.eps}{\special{language
"Scientific Word";type "GRAPHIC";maintain-aspect-ratio TRUE;display
"ICON";valid_file "F";width 11.2467cm;height 8.468cm;depth
0pt;original-width 8.003in;original-height 6.0105in;cropleft "0";croptop
"1.0001";cropright "0.9996";cropbottom "0";filename
'BandStop_Cheb_dB_Res.eps';file-properties "XNPEU";}}\FRAME{ftbphFU}{%
11.2489cm}{8.468cm}{0pt}{\Qcb{Deviation between direct and backward
calculations of the band-stop Chebyshev-1 magnitude on a 32-bits processor.}%
}{\Qlb{BandStop_Cheb1_Error}}{bandstop_cheb_error.eps}{\special{language
"Scientific Word";type "GRAPHIC";maintain-aspect-ratio TRUE;display
"ICON";valid_file "F";width 11.2489cm;height 8.468cm;depth
0pt;original-width 8.003in;original-height 6.0105in;cropleft "0";croptop
"1";cropright "0.9998";cropbottom "0";filename
'BandStop_Cheb_Error.eps';file-properties "XNPEU";}} 
%TCIMACRO{\TeXButton{clearpage}{\clearpage}}%
%BeginExpansion
\clearpage%
%EndExpansion
\end{center}

Figs. \ref{Low_Butt_Res} and \ref{Low_Butt_dB_Res} show the magnitude
response of a low-pass, Butterworth filter, while Fig. \ref{Low_Butt_Error}
shows the deviation between the two methods. The filter order in this case
is $N=10$.

Figs. \ref{Low_Cheb1_Res}, \ref{Low_Cheb1_dB_Res} and \ref{Low_Cheb1_Error}
show the results of the same input data but for a Chebyshev-1 filter. As we
see from Fig. \ref{Low_Cheb1_Res}, we manage with $N=6$.

Figs. \ref{High_Butt_Res} and \ref{High_Butt_dB_Res} show the magnitude
response of a high-pass, Butterworth filter, while Fig \ref{High_Butt_Error}%
, again, shows the deviation between the two methods. The filter order is $%
N=16$.

Figs. \ref{High_Cheb1_Res}, \ref{High_Cheb1_dB_Res} and \ref%
{High_Cheb1_Error} show the results of the same input data but for a
Chebyshev-1 filter. In this case we are able to reduce the filter order to $%
N=8$.

Figs. \ref{BandPass_Butt_Res} and \ref{Bandpass_Butt_dB_Res} show the
magnitude response of a band-pass, Butterworth filter, while Fig. \ref%
{BandPass_Butt_Error} shows the deviation between the two methods. The
filter order in this case is $N=20$.

Figs. \ref{BandPass_Cheb1_Res}, \ref{BandPass_Cheb1_dB_Res} and \ref%
{BandPass_Cheb1_Error} show the results of the same input data but for a
Chebyshev-1 filter. Filter order in this case is reduced to $N=12$.

Figs. \ref{BandStop_Butt_Res} and \ref{BandStop_Butt_dB_Res} show the
magnitude response of a band-stop, Butterworth filter, while Fig. \ref%
{BandStop_Butt_Error} shows the deviation between the two methods. The
filter order in this case is $N=20$.

Figs. \ref{BandStop_Cheb1_Res}, \ref{BandStop_Cheb1_dB_Res} and \ref%
{BandStop_Cheb1_Error} show the results of the same input data but for a
Chebyshev-1 filter. Filter order in this case is reduced to $N=12$.

The input and design parameters of Butterworth and Chebyshev-1 filters are
gathered in Tabs. \ref{ButterTable} and \ref{Cheb1Table}.

\begin{center}
%TCIMACRO{\TeXButton{B}{\begin{table}[tbp] \centering}}%
%BeginExpansion
\begin{table}[tbp] \centering%
%EndExpansion
\begin{tabular}{|c|c|c|c|c|c|}
\hline
& $f_{\lim }$ & $A_{\lim }\left( dB\right) $ & Word Length & $N$ & $f_{c}$
\\ \hline
Low-pass & $0.2,\,0.25$ & $1,\,20$ & $16$ & $10$ & $0.2103$ \\ \hline
High-pass & $0.3,\,0.35$ & $40,\,1$ & $32$ & $16$ & $0.3063$ \\ \hline
Band-pass & $0.2,\,0.25,\,0.35,\,0.4$ & $20,\,1,\,1,\,30$ & $32$ & $20$ & $%
0.2393,\,0.3585$ \\ \hline
Band-stop & $0.2,\,0.3,\,0.4,\,0.45$ & $1,\,20,\,35,\,1$ & $32$ & $16$ & $%
0.2129,\,0.4457$ \\ \hline
\end{tabular}
\caption{Input and design parameters for the Butterworth
filter.\label{ButterTable}} 
%TCIMACRO{\TeXButton{E}{\end{table}}}%
%BeginExpansion
\end{table}%
%EndExpansion

%TCIMACRO{\TeXButton{B}{\begin{table}[tbp] \centering}}%
%BeginExpansion
\begin{table}[tbp] \centering%
%EndExpansion
\begin{tabular}{|c|c|c|c|c|c|}
\hline
& $f_{\lim }$ & $A_{\lim }\left( dB\right) $ & Word Length & $N$ & $\epsilon 
$ \\ \hline
Low-pass & $0.2,\,0.25$ & $1,\,20$ & $16$ & $6$ & $0.5088$ \\ \hline
High-pass & $0.3,\,0.35$ & $40,\,1$ & $32$ & $8$ & $0.5088$ \\ \hline
Band-pass & $0.2,\,0.25,\,0.35,\,0.4$ & $20,\,1,\,1,\,30$ & $32$ & $12$ & $%
0.5088$ \\ \hline
Band-stop & $0.2,\,0.3,\,0.4,\,0.45$ & $1,\,20,\,35,\,1$ & $32$ & $12$ & $%
0.5088$ \\ \hline
\end{tabular}
\caption{Input and design parameters for the Chebyshev-1 filter.
\label{Cheb1Table}} 
%TCIMACRO{\TeXButton{E}{\end{table}}}%
%BeginExpansion
\end{table}%
%EndExpansion
\end{center}

\section{SUMMARY AND CONCLUSION}

We began this report by describing how a digital low-pass, high-pass,
band-pass and band-stop filters might be specified. Then we introduced
Butterworth and Chebyshev-1 analog filters and explained how their
parameters could be deduced from the design specifications. By applying the
Bilinear Method we transformed this prototype to a digital filter. The
results were of course the decimal filter coefficients. We called this
approach the Direct calculations, scaled and saved both decimal and integer
coefficients. The Backward method, then applied the integer coefficients to
construct the impulse-response and $\mathcal{F}$-transform it to obtain the
frequency-response. The magnitude-responses of both methods were then
calculated. In the last section, we compared the results and plotted the
deviation between the two methods for some different set of input
parameters. We may conclude the following:

\begin{itemize}
\item All digital filters implemented here are unconditionally stable and
causal.

\item All filters are implemented by transforming a corresponding analog
low-pass filter to a digital low-pass using the Bilinear transformation,
then applying frequency mapping to convert this prototype to a low-pass,
high-pass, band-pass and band-stop filters with different cut-off
frequencies.

\item The coefficients of a band-pass filter may also be calculated by
cascading a high-pass and a low-pass filter. We have not used this approach
because of two reasons. First, it is not consistent with the high-pass and
band-stop implementation. Second, it gives less satisfactory results.

\item While a low-pass Butterworth filter has a flat amplitude in both ends,
the frequency of oscillation in the pass-band side of a low-pass Chebyshev-1
grows as we approach the $\omega _{1}$, just as their analog counter-parts.

\item Given the same set of specifications, Chebyshev-1 filters are
advantageous because of lower filter order. The price is, some oscillation
in the pass-band, as well as a little more complexity of implementing them.

\item For a low-pass Butterworth filter, the deviation between the two
methods is flat in both ends and reaches its maximum value at $\Omega _{c}$.
For a low-pass Chebyshev-1 filter, the deviation is flat in the pass-band,
oscillatory in the stop-band and reaches its maximum value at $\Omega _{1}$.

\item For the low-pass and high-pass filters, the results are quite
satisfactory, and the deviation could actually be seen just around the
critical frequencies and in the $dB$-plot for extremely low values of the
magnitude responses.

\item For the band-pass and band-stop cases we get a little more deviation
between the two methods, although still quite satisfactory. This is
specially the case if we require steep transitions near the edges of the
frequency interval. It comes from the fact that these cases have a filter
order twice the corresponding low-pass filter, resulting very small
coefficients and more numeric problems.

\item As we could imagine, the deviation decreases as we increase the
impulse-response length and the processor capabilities.
\end{itemize}

\section{REFERENCES}

\begin{thebibliography}{9}
\bibitem{Mitra} Sanjit K. Mitra and James F. Kaiser: Handbook for Digital
Signal Processing, 1993, John Wiley \& Sons, NY, Chapter 5, ISBN\ \ \
0-471-61995-7.

\bibitem{Oppenheim} Alan V. Oppenheim and Roland W. Schafer, Discrete-Time
Signal Processing, 1989, Prentice-Hall International, Inc., Chapters 2, 3, 4
and 5, 7 and 8, ISBN\ \ \ 0-13-216771-9.

\bibitem{Smith} The Scientist and Engineer's Guide to Digital Signal
Processing, 1997, Steven W. Smith, California Technical Publishing, Chapters
20 and 33, ISBN\ \ \ 0-9660176-3-3

%TCIMACRO{\TeXButton{clearpage}{\clearpage}}%
%BeginExpansion
\clearpage%
%EndExpansion
\end{thebibliography}

\end{document}
